\documentclass[journal]{IEEEtran}

% ============================================================================
% PAKETLER
% ============================================================================
\usepackage[utf8]{inputenc}
\usepackage[T1]{fontenc}
\usepackage[turkish,shorthands=off]{babel}
\usepackage{cite}
\usepackage{amsmath,amssymb,amsfonts}
\usepackage{algorithmic}
\usepackage{graphicx}
\usepackage{textcomp}
\usepackage{xcolor}
\usepackage{booktabs}
\usepackage{multirow}
\usepackage{subcaption}
\usepackage{array}
\usepackage{balance}
\usepackage{hyperref}

% Ozel komutlar
\newcommand{\eps}{\epsilon}
\newcommand{\Linf}{L_\infty}
\renewcommand{\IEEEkeywordsname}{Anahtar Kelimeler}

% NOT: Bu surum, TR Dizin dergi gonderimi icin calisma taslagi formatidir
% (IEEEtran). Secilen derginin kendi sablonuna aktarim gonderim asamasinda
% yapilacaktir; cogu TR Dizin dergisi Word/serbest bicim kabul eder.
% Ingilizce ozgun metin: paper/manuscript/ (birebir ayni bilimsel icerik).

% ============================================================================
% BELGE BILGISI
% ============================================================================
\title{Ölçüm Protokolü Sonucu Belirler: CNN ile Görü Dönüştürücüsü Çekişmeli Gürbüzlük Karşılaştırmalarının Yöntembilimsel Bir İncelemesi\\
{\normalsize\itshape The Measurement Protocol Decides the Conclusion: A Methodological Study of CNN and Vision Transformer Adversarial Robustness Comparisons}}

\author{Çağrı~Şahin ve~Özal~Yıldırım%
\thanks{Çağrı Şahin, Muş Alparslan Üniversitesi Yazılım Mühendisliği Bölümü'ndedir (49250 Muş, Türkiye) ve ayrıca Fırat Üniversitesi Yazılım Mühendisliği Bölümü'nde doktora öğrencisidir (23119 Elazığ, Türkiye).}%
\thanks{Özal Yıldırım, Fırat Üniversitesi Yazılım Mühendisliği Bölümü'ndedir, 23119 Elazığ, Türkiye.}%
\thanks{Sorumlu yazar: Çağrı Şahin (e-posta: c.sahin@alparslan.edu.tr).}}

% ============================================================================
% BELGE
% ============================================================================
\begin{document}

\maketitle

% ----------------------------------------------------------------------------
% OZ
% ----------------------------------------------------------------------------
\begin{abstract}
Mimari aileler arasındaki çekişmeli gürbüzlük karşılaştırmaları literatürde çelişen sonuçlar bildirmekte ve bu uyuşmazlık genellikle modellere, eğitim tariflerine veya veri kümelerine atfedilmektedir. Bu çalışma, uyuşmazlığın önemli bir bölümünün ölçüm protokolünün kendisinden geldiğini göstermektedir. CIFAR-10 üzerinde çekişmeli eğitilmiş ResNet-18 ve ViT-Tiny çiftleri kullanılarak (mimari başına üç bağımsız eğitim koşusu, örnek bazında kayıtla tam test kümesi değerlendirmesi ve her eğitim aşamasından önce ayrılmış sabit bir doğrulama bölmesi) yalnızca gürbüzlüğün ve transferin nasıl ölçüldüğü değiştirilmiştir. Dört yerleşik koşullama protokolü boyunca ölçülen CNN$\rightarrow$ViT transfer asimetrisi 4,4 ile 14,6 puan arasında değişmekte (3,3 katlık bir yayılım) ve protokol seçimi bir eşdeğerlik testinin kabul mü ret mi edileceğine karar vermektedir; buna karşılık etkinin yönü kararlıdır. Tasarımın tamamı CIFAR-100 üzerinde yinelendiğinde tohum başına yayılım 10,5 puandan 13,6 puana çıkmakta ve işaret on iki ölçümün tamamında korunmaktadır; dolayısıyla etki tek bir veri kümesine özgü değildir. Eğitim yörüngesini sabit tutup yalnızca çevrimdışı kontrol noktası seçimi kuralını değiştiren ek bir ablasyon, raporlanan gürbüz doğruluğu 1,6 ile 2,9 puan arasında oynatmaktadır; bu, neredeyse hiç raporlanmayan dördüncü bir varyans kaynağıdır. Üçüncü bir veri kümesi bunu keskinleştirmektedir. Aynı tarifin iki mimariyi temiz doğrulukta yalnızca 1,85 puan ayırdığı SVHN'de dört protokolün ikisi CNN üstünlüğü, ikisi ViT üstünlüğü bildirmekte ve birincil protokol istatistiksel eşdeğerliği kabul etmektedir: modeller eşitlendiğinde ölçüm tercihi, raporlanan asimetrinin yalnızca büyüklüğüne değil işaretine de karar vermektedir. WideResNet-28-10 referanslı $3\times3$ transfer matrisi, koşulsuz oranların koşullu oranlardan sapmasının hedefin temiz hatasıyla yalnızca ilişkili olmadığını, ona cebirsel olarak eşit olduğunu göstermektedir: temiz girdide yanlış sınıflandırılan örnekler saldırı altında da yanlış kaldığı için (36 yönde ölçüm 0,989-1,000), ham = temiz hata + koşullu oran $\times$ (1 - temiz hata) özdeşliği 0,41 puandan küçük bir hatayla tutmaktadır. Ham transfer oranları bu yüzden onları raporlayan her çalışmada yapısı gereği hedef zayıflığıyla kirlenmektedir. Koşullu ayrıştırma, CNN üstünlüğünü temiz doğruluk ve koşullu duyarlılık çarpanlarına ayırmakta; mutlak sıralama değişmediği hâlde iki çarpan arasındaki pay eğitim koşuları arasında değişebilmektedir. Son olarak, bu literatürde yaygın üç iddia doğrudan ölçüm altında tekrarlanamamıştır: girdi gradyanlarında mekânsal lokalite farkı yoktur, CLS dikkat entropisi $n=1000$'de saldırı altında değişmemektedir ve ViT'e özgü bir transfer saldırısı bu rejimde transfer kazancı sağlamamaktadır. Ölçekten bağımsız gradyan seyrekliği farkları ile çekişmeli hasarın mimariye özgü derinlik profilleri ise ayakta kalmaktadır. Çalışma, mimariler arası gürbüzlük çalışmaları için raporlama gereklilikleriyle sonlanmaktadır.
\end{abstract}

\begin{IEEEkeywords}
Çekişmeli gürbüzlük, değerlendirme yöntembilimi, derin öğrenme güvenliği, evrişimli sinir ağları, görü dönüştürücüsü, tekrarlanabilirlik, transfer saldırıları
\end{IEEEkeywords}

% ----------------------------------------------------------------------------
% BOLUMLER
% ----------------------------------------------------------------------------
% ============================================================================
% BOLUM 1: GIRIS
% ============================================================================
\section{Giriş}
\label{sec:introduction}

Derin sinir ağları, küçük bir norm bütçesiyle sınırlı kalırken yüksek güvenle yanlış sınıflandırmaya yol açan, özenle üretilmiş pertürbasyonlar olan çekişmeli örneklere (adversarial examples) karşı savunmasızdır~\cite{szegedy2014intriguing, goodfellow2015explaining}. Görü Dönüştürücüleri (Vision Transformer, ViT)~\cite{dosovitskiy2021image, vaswani2017attention} evrişimli sinir ağlarının (CNN)~\cite{he2016deep} yanında baskın bir mimari aile hâline gelince doğal bir soru ortaya çıkmıştır: hangi aile daha gürbüzdür ve neden?

Literatür bu soruya tutarsız yanıtlar vermektedir. Bai vd.~\cite{bai2021transformers} eşleşmiş eğitim tarifleri altında dönüştürücülerin CNN'leri geçmediğini, Benz vd.~\cite{benz2021adversarial} ise standart eğitilmiş modellerde ViT'lerin daha gürbüz olduğunu bildirmektedir; mimariler arası transfer çalışmaları da saldırıların CNN'den ViT'e mi yoksa tersine mi daha kolay taşındığı konusunda uyuşmamaktadır~\cite{mahmood2021robustness, ravikumar2023trend, naseer2022improving}. Bu uyuşmazlıklar genellikle model, eğitim tarifi veya veri kümesi farklarına atfedilir ve standart tepki, eşleşmiş bir tarif altında karşılaştırma yapmaktır.

Bu makale farklı bir gözlemde bulunmaktadır. Modeller, veri, tehdit modeli ve saldırı bütçesinin tamamı sabit tutulduğunda bile sonuç, nadiren raporlanan bir seçime bağlı kalmaktadır: \emph{ölçüm protokolü}. Transfer çalışmaları paydaya hangi örneklerin gireceğine karar vermek zorundadır (tüm test örnekleri mi, hedefin doğru sınıflandırdıkları mı, her iki modelin de doğru sınıflandırdıkları mı, yoksa kaynak saldırının başarılı olduğu örnekler mi); gürbüzlük karşılaştırmaları da mutlak bir skor mu raporlayacaklarına yoksa onu temiz doğruluk ve koşullu duyarlılık bileşenlerine mi ayıracaklarına karar vermek zorundadır. Bu seçimler genellikle raporlama ayrıntısı sayılır. Biz bunların sonucun parçası olduğunu gösteriyoruz.

% ----------------------------------------------------------------------------
% Yaklasim
% ----------------------------------------------------------------------------
Kurulumumuz, ölçüm sorusunu yalıtabilmek için bilinçli olarak dardır. CIFAR-10~\cite{krizhevsky2009learning} üzerinde ResNet-18 ile ViT-Tiny'yi eşleşmiş bir amaç fonksiyonu ve tehdit modeliyle çekişmeli eğitiyoruz; mimari başına üç bağımsız eğitim koşusu yapılmakta ve \emph{her} eğitim aşamasından önce ayrılan sabit bir doğrulama bölmesi kullanılmaktadır, dolayısıyla hiçbir model gördüğü veriyle seçilmemektedir. Tüm değerlendirmeler, örnek bazında sonuç kaydıyla tam test kümesinde koşmaktadır; bu, eşleştirilmiş istatistikleri mümkün kılar: aynı örnekler, saldırıyı yeniden koşmadan farklı protokoller altında yeniden ağırlıklandırılabilir. Ardından modelleri değil protokolü değiştiriyor ve sonucun ne kadar hareket ettiğini ölçüyoruz.

Mimari karşılaştırma makalelerinde tekrar eden üç iddiayı da yorum olarak değil, ölçülebilir hipotez olarak ele alıyoruz: CNN girdi gradyanlarının mekânsal olarak daha lokalize olduğu, çekişmeli pertürbasyonun ViT dikkatini (attention) bozduğu ve ViT'e özgü transfer saldırılarının mimariler arası transferi iyileştirdiği. Her biri, amaca özel bir ölçümle doğrudan test edilmektedir.

% ----------------------------------------------------------------------------
% Katkilar
% ----------------------------------------------------------------------------
Katkılarımız şunlardır.

\begin{enumerate}
\item \textbf{Protokol duyarlılığının nicelenmesi.} Aynı modeller ve aynı veri üzerinde, ölçülen mimariler arası transfer asimetrisi dört yerleşik koşullama protokolü boyunca $+4{,}4$ ile $+14{,}6$ puan arasında değişmektedir; bu 3,3 katlık bir yayılımdır; buna karşılık her iki mimariyi sıfırdan yeniden eğitmek aynı niceliği $1{,}5$ puandan az oynatmaktadır. Protokol, bir eşdeğerlik testinin kabul mü ret mi edileceğine karar vermektedir; asimetrinin yönü ise büyüklüğünün aksine kararlıdır.

\item \textbf{Etkinin ikinci bir veri kümesinde yinelenmesi ve dördüncü bir varyans kaynağı.} Tasarımın tamamı CIFAR-100 üzerinde tekrarlandığında tohum başına ortalama protokol yayılımı $13{,}58\pm1{,}71$ puana çıkmakta ve işaret on iki ölçümün tamamında korunmaktadır. Ayrıca bir eğitim yörüngesi sabit tutulup yalnızca çevrimdışı kontrol noktası seçimi kuralı değiştirildiğinde raporlanan PGD-10 doğruluğu $1{,}58$--$2{,}85$ puan oynamaktadır; bu, seçimi tohumlar, saldırılar ve protokollerin yanında raporlanması gereken bir varyans kaynağı hâline getirmektedir.

\item \textbf{Koşulsuz transfer oranlarındaki karıştırıcının ölçülmesi.} WideResNet-28-10 referansının eklendiği $3\times3$ transfer matrisinde, ham oranların koşullu oranlardan sapması ampirik bir korelasyon değil bir özdeşliktir: hedefin zaten yanlış sınıflandırdığı örnekler pertürbasyonla neredeyse hiç kurtulmadığı için ($36$ yönde $0{,}989$--$1{,}000$), ham oran, temiz hata artı bir eksi o hatayla ölçeklenmiş koşullu orana $0{,}41$ puandan küçük bir farkla eşittir. Ayrıca gelen transferin kaynak mimarisini değil hedefin kendi kırılganlığını izlediğini buluyoruz (üç hedef üzerinde $r=0{,}986$); bu, asimetriyi güçlü CNN saldırılarının değil, daha zayıf hedefin bir özelliği olarak yeniden çerçevelemektedir.

\item \textbf{Koşullu ayrıştırma ve atfının kararsızlığına dair kanıt.} Mutlak gürbüzlük, temiz doğruluk ile koşullu duyarlılığın çarpımına birebir ayrışır. İkisini birden raporlamak önemlidir; çünkü mutlak sıralama değişmediği hâlde, iki çarpan arasındaki pay önceki tek koşuluk kontrol noktalarımız ile üç yeni koşu arasında değişmiştir.

\item \textbf{Üç negatif sonuç.} Ölçekten bağımsız seyreklik farkları mekânsal lokalite farkına karşılık gelmemektedir (dört mekânsal ölçütten üçü sıfır sonuç); CLS dikkat entropisi $n=1000$'de saldırı altında ölçülebilir biçimde değişmemektedir; ViT'e özgü bir transfer saldırısı olan Jeton Gradyanı Düzenlileştirme (TGR)~\cite{zhang2023tgr} bu rejimde transfer kazancı sağlamamakta, eşleşmiş bütçede düz MI-FGSM'ye yenilmektedir.

\item \textbf{Ayakta kalanlar.} Gradyan seyrekliği farkları bağımsız kontrol noktası kümelerinde tekrarlanmakta, mimariden değil çekişmeli eğitimden kaynaklanmakta ve yalnızca mutlak kosinüste var olan bir hizalanma farkıyla birlikte gelmektedir. Çekişmeli hasar iki mimaride nitel olarak farklı derinliklerde birikmektedir ve ölçülen kayma profili, hesaplamada kullanılan jeton (token) agregasyonuna bağlıdır.
\end{enumerate}

Çalışmayı, bu ölçümlerden çıkan kısa bir raporlama gereklilikleri listesiyle kapatıyoruz.

% ----------------------------------------------------------------------------
% Makale Duzeni
% ----------------------------------------------------------------------------
Bölüm~\ref{sec:related_work} çekişmeli değerlendirme pratiğini ve mimari karşılaştırma çalışmalarını gözden geçirmektedir. Bölüm~\ref{sec:methodology} eğitim protokolünü, koşullama protokollerini ve ölçümleri tanımlamaktadır. Bölüm~\ref{sec:experiments} sonuçları raporlamakta, Bölüm~\ref{sec:discussion} çıkarımları ve sınırlılıkları tartışmakta, Bölüm~\ref{sec:conclusion} çalışmayı sonuçlandırmaktadır.

% ============================================================================
% BOLUM 2: ILGILI CALISMALAR
% ============================================================================
\section{İlgili Çalışmalar}
\label{sec:related_work}

% ----------------------------------------------------------------------------
% 2.1 Cekismeli Saldirilar
% ----------------------------------------------------------------------------
\subsection{Çekişmeli Saldırılar}
\label{subsec:attacks}

Çekişmeli saldırılar, yanlış sınıflandırmaya yol açan en küçük pertürbasyonları bulmayı amaçlar. Hızlı Gradyan İşaret Yöntemi (FGSM)~\cite{goodfellow2015explaining}, kayıp gradyanının işaretini izleyerek pertürbasyonu tek adımda üretir; pertürbasyon büyüklüğü $\Linf$ normu altında sınırlıdır. İzdüşümlü Gradyan İnişi (PGD)~\cite{madry2018towards} FGSM'yi yinelemeli eniyilemeyle genişletir: birden çok gradyan-işaret adımı atar ve her adımdan sonra $\eps$-topuna geri izdüşüm yapar (biçimsel tanımlar Bölüm~\ref{subsec:attack_methods}).

Carlini--Wagner (C\&W) saldırısı~\cite{carlini2017towards} çekişmeli örnek üretimini bir eniyileme problemi olarak kurar; genellikle daha yüksek başarı oranlarına ulaşır ancak hesaplama maliyeti artar. Daha yakın tarihli AutoAttack~\cite{croce2020reliable}, güvenilir bir gürbüzlük değerlendirmesi sağlamak için birden çok saldırı stratejisini (APGD-CE, APGD-T, FAB-T, Square) birleştirir ve kıyaslamada fiilî standart hâline gelmiştir.

% ----------------------------------------------------------------------------
% 2.2 Savunma Mekanizmalari
% ----------------------------------------------------------------------------
\subsection{Savunma Mekanizmaları}
\label{subsec:defenses}

Çekişmeli eğitim (adversarial training, AT)~\cite{madry2018towards} en etkili savunma olmayı sürdürmektedir; modeli her partide üretilen çekişmeli örneklerle eğitir:
\begin{equation}
    \min_\theta \mathbb{E}_{(x,y) \sim \mathcal{D}} \left[ \max_{\|\delta\|_\infty \leq \eps} \mathcal{L}(f_\theta(x + \delta), y) \right]
\end{equation}

TRADES~\cite{zhang2019theoretically}, doğal doğruluk ile gürbüzlük arasında kuramsal temelli bir ödünleşim sunar:
\begin{equation}
    \mathcal{L}_{\text{TRADES}} = \mathcal{L}_{CE}(f(x), y) + \beta \cdot \text{KL}(f(x) \| f(x_{adv}))
\end{equation}
burada $\beta$ ödünleşim parametresidir.

MART~\cite{wang2020improving}, yanlış sınıflandırılan örnekleri eğitim sırasında vurgulayarak ele alır. Diğer yaklaşımlar arasında girdi dönüşümü~\cite{guo2018countering}, rastgeleleştirilmiş yumuşatma yoluyla sertifikalı savunmalar~\cite{cohen2019certified} ve topluluk yöntemleri~\cite{tramer2018ensemble} bulunur. Sertifikalı gürbüzlükteki güncel gelişmeler arasında, çok adımlı savunmalarla sertifika yarıçaplarını iyileştiren Uyarlanabilir Rastgeleleştirilmiş Yumuşatma~\cite{lyu2024adaptive} sayılabilir.

% ----------------------------------------------------------------------------
% 2.3 CNN Gurbuzlugu
% ----------------------------------------------------------------------------
\subsection{CNN Gürbüzlüğü}
\label{subsec:cnn_robustness}

CNN mimarilerinin çekişmeli gürbüzlüğü kapsamlı biçimde incelenmiştir. Madry vd.~\cite{madry2018towards}, PGD saldırılarıyla çekişmeli eğitimin güçlü deneysel gürbüzlük sağladığını göstermiştir. RobustBench lider tablosu~\cite{croce2021robustbench} en güncel CNN gürbüzlüğünü izlemektedir; CIFAR-10'da en yüksek gürbüz doğruluğa WideResNet türevleri ulaşmaktadır.

Model kapasitesinin ulaşılabilir gürbüzlükle ilişkili olduğu gösterilmiştir~\cite{madry2018towards}; daha geniş ve derin ağlar daha iyi gürbüz doğruluk sergilemektedir. Mimariye özgü bulgular arasında atlama bağlantılarının çekişmeli transfer edilebilirlikteki rolü~\cite{wu2020skip} ve çekişmeli eğitimin büyük ölçekte yol açtığı ayrık parti-normalizasyonu istatistikleri~\cite{xie2020intriguing} sayılabilir; ilişkili olarak, yardımcı parti-normalizasyonu dalları çekişmeli örneklerin temiz doğruluğu artırmasına imkân vermektedir~\cite{xie2020adversarial}. Çekişmeli eğitimin girdi-gradyan yapısının kendisini de yeniden biçimlendirdiği, daha algısal hizalı~\cite{tsipras2019robustness} ve daha seyrek~\cite{chalasani2020concise} gradyanlar ürettiği bilinmektedir.

% ----------------------------------------------------------------------------
% 2.4 ViT Gurbuzlugu
% ----------------------------------------------------------------------------
\subsection{Görü Dönüştürücüsü Gürbüzlüğü}
\label{subsec:vit_robustness}

Vaswani vd.~\cite{vaswani2017attention} tarafından tanıtılan öz-dikkat mekanizması üzerine kurulan Görü Dönüştürücüleri~\cite{dosovitskiy2021image}, görüntüleri yama dizileri olarak işler; iç temsilleri CNN'lerden katman katman ölçülebilir biçimde ayrışır~\cite{raghu2021vision}. İlk çalışmalar, küresel alıcı alan sayesinde ViT'lerin bazı içkin gürbüzlük özellikleri sergilediğini öne sürmüştür~\cite{bhojanapalli2021understanding}. Bai vd.~\cite{bai2021transformers}, dönüştürücülerin CNN'lerden daha gürbüz olup olmadığını sistematik olarak incelemiş; adil eğitim tarifleri altında çekişmeli gürbüzlükte üstün olmadıklarını, buna karşılık öz-dikkate atfedilen dağılım kayması gürbüzlüğünde öne geçtiklerini bulmuştur. Benz vd.~\cite{benz2021adversarial}, MLP-Mixer'i de içeren erken bir doğrudan karşılaştırma raporlamış ve farkları frekans tercihlerine bağlamıştır. Sonraki çalışmalar ViT'lerin özellikle gradyan tabanlı saldırılara karşı savunmasız kaldığını~\cite{mahmood2021robustness}, yama hedefli saldırıların dikkat mekanizmasını doğrudan istismar ettiğini~\cite{fu2022patchfool} ve dikkatın yama pertürbasyonları altında ölçülebilir biçimde kaydığını~\cite{gu2022vision} ortaya koymuştur.

Shao vd.~\cite{shao2022adversarial}, çekişmeli eğitimin ViT gürbüzlüğünü artırabildiğini ve ViT özniteliklerinin yüksek frekanslı içeriğe daha az dayandığını göstermiştir; eğitim kararsızlığı kaynaklı güçlükler ise sürmektedir. ViT'lerin yaygın bozulmalara ve örtmelere karşı gürbüzlüğü de betimlenmiştir~\cite{paul2022vision}. Daha yakın zamanda Jain vd.~\cite{jain2024towards}, ViT'e özgü eğitim güçlüklerini ele alan AAS-AT yöntemini önermiş; difüzyonla artırılmış çekişmeli eğitimle ViT'lerin küçük ölçekli veri kümelerinde WideResNet'leri geçebildiği bildirilmiştir~\cite{wu2025vision}. Zhu vd.~\cite{zhu2024enhancing}, gradyan normalizasyonu ölçekleme ve yüksek frekans uyarlamasıyla ViT'lere yönelik transfer edilebilir saldırıları güçlendirmiştir. Bu ilerlemelere karşın, birleşik değerlendirme çerçeveleri altında kapsamlı CNN--ViT davranışsal karşılaştırmaları sınırlı kalmaktadır.

% ----------------------------------------------------------------------------
% 2.5 Transfer Saldirilari
% ----------------------------------------------------------------------------
\subsection{Transfer Saldırıları}
\label{subsec:transfer_related}

Çekişmeli örnekler modeller arasında sıklıkla transfer olur ve bu, kara kutu saldırılarını mümkün kılar~\cite{papernot2016transferability}. Transfer edilebilirlik; model mimarisi, eğitim verisi ve saldırı yöntemi gibi etkenlere bağlıdır~\cite{demontis2019adversarial}. Temel transfer çalışmaları transfer edilebilirliği zaten doğru sınıflandırılan örnekler üzerinde değerlendirmekte~\cite{liu2017delving}, standart saldırı makaleleri de başarı oranlarını aynı gerekçeyle koşullamaktadır~\cite{dong2018boosting}; güncel değerlendirme kılavuzları bu tür protokol seçimlerini açık hâle getirmektedir~\cite{zhao2025revisiting}. Topluluk tasarımında Ravikumar vd.~\cite{ravikumar2023trend}, transfer ölçümünü tam olarak modeller arasındaki doğruluk farklarını hesaba katmak için her iki modelin de doğru sınıflandırdığı örneklerle sınırlamaktadır --- koşullu metriğimizin kontrol ettiği karıştırıcının aynısı.

CNN'ler ile ViT'ler arasındaki mimariler arası transfer artan ilgi görmektedir. Mahmood vd.~\cite{mahmood2021robustness}, her-ikisi-doğru protokolüyle mimariler arası transferi ölçmüş ve dikkat çekici ölçüde düşük CNN$\leftrightarrow$ViT transferi bulmuştur; Naseer vd.~\cite{naseer2022improving}, ViT'lerde üretilen çekişmeli örneklerin geleneksel saldırılarla zayıf transfer ettiğini göstermiş ve ViT kaynaklı transferi iyileştirmek için öz-topluluk ve jeton arıtma önermiştir; Wei vd.~\cite{wei2022towards} dikkat gradyanlarını atlayarak ve yama bazlı seyreltmeyle ViT kaynaklı transferi geliştirmiştir. Bunların üzerine Jeton Gradyanı Düzenlileştirme (TGR)~\cite{zhang2023tgr}, ViT'in dikkat bileşenlerini kullanarak transfer edilebilir örnekler üretmekte; Momentum Bütünleşik Gradyanlar (MIG)~\cite{ma2023transferable} hem ViT'lerde hem CNN'lerde gelişmiş transfer sağlamaktadır. Chen vd.~\cite{chen2023understanding}, mimariler arası transferi etkileyen etkenleri sistematik olarak incelemekte ve mimari tümevarım önyargılarının (CNN'lerin yüksek, ViT'lerin düşük frekanslı bilgiyi tercih etmesi) transfer başarısını önemli ölçüde etkilediğini bulmaktadır. Bu transfer örüntülerini anlamak, gürbüz topluluk sistemleri geliştirmek ve gerçek dünya saldırı senaryolarını değerlendirmek için önemlidir.

% ----------------------------------------------------------------------------
% 2.6 Bosluk ve Konumumuz
% ----------------------------------------------------------------------------
\subsection{Araştırma Boşluğu}
\label{subsec:gap}

Tekil çalışmalar CNN ve ViT gürbüzlüğünü ayrı ayrı incelemiş ve adil olmayan CNN--dönüştürücü karşılaştırmalarının tuzaklarına daha önce dikkat çekilmiş olsa da~\cite{pinto2022impartial, bai2021transformers}, mevcut karşılaştırmalı çalışmalar gürbüzlük skorlarını çoğunlukla tek başına, üstelik tutarsız pertürbasyon bütçeleri, saldırı yapılandırmaları veya model ölçekleriyle raporlamakta; karşılaştırmayı davranışsal analizle eşleştirmemektedir. Yakın tarihli bir çalışma ViT ve CNN çekişmeli gürbüzlüğünü doğrudan karşılaştırmaktadır~\cite{ali2024adversarial}; çalışmamız, tam test kümesinde AutoAttack ve eşleştirilmiş istatistiklerle değerlendirme yapması, transfer ölçümünü temiz doğruluk farklarına karşı kontrol etmesi ve karşılaştırmaya gradyan ile öznitelik düzeyinde davranışsal analizler eklemesiyle ayrışmaktadır.

Çalışmamız ölçüm sorusunu doğrudan ele almakta ve onu incelemenin nesnesi hâline getirmektedir. Transfer değerlendirmesini doğru sınıflandırılan örneklerle koşullamak yerleşik pratiktir~\cite{liu2017delving, dong2018boosting, mahmood2021robustness, ravikumar2023trend} ve tekil makaleler bir protokol seçip devam etmektedir. Eksik olan, bu seçimin ne kadar önemli olduğuna dair kontrollü bir tahmindir. Bunu; modelleri, veriyi, tehdit modelini ve saldırı bütçesini sabit tutup yalnızca protokolü değiştirerek sağlıyor ve ölçülen asimetride 3,3 katlık bir yayılım buluyoruz --- bağımsız eğitim koşuları arasındaki standart sapmanın yaklaşık yirmi katı. Koşulsuz oranlardaki karıştırıcıyı üçüncü bir mimariyle de niceliyoruz (hedefin temiz hatasına karşı $r=0{,}997$); bu, bilinen nitel bir uyarıyı yayımlanmış sonuçlara karşı sınanabilir bir sayıya dönüştürmektedir.

Bu literatürde tekrar eden ve tasarımımızın varsaymak yerine test etmeye imkân verdiği üç iddia daha vardır: CNN gradyanlarının mekânsal olarak daha lokalize olduğu, çekişmeli pertürbasyonun ViT dikkatini bozduğu ve dönüştürücüye özgü transfer saldırılarının~\cite{zhang2023tgr, naseer2022improving} mimariler arası transferi iyileştirdiği. Her birini doğrudan ölçüyoruz ve hiçbiri bizim kurulumumuzda geçerli çıkmıyor --- bu sonuçları raporluyoruz, çünkü protokole duyarlı bir karşılaştırmanın kaldıramayacağı şey tam da bu türden test edilmemiş bir varsayımdır.

% ============================================================================
% BOLUM 3: YONTEM
% ============================================================================
\section{Yöntem}
\label{sec:methodology}

% ----------------------------------------------------------------------------
% 3.1 Tehdit Modeli
% ----------------------------------------------------------------------------
\subsection{Tehdit Modeli}
\label{subsec:threat_model}

Saldırganın girdi görüntülerini sınırlı bir bölge içinde pertürbe edebildiği standart $\Linf$ tehdit modelini ele alıyoruz:
\begin{equation}
    \mathcal{B}_\eps(x) = \{x' : \|x' - x\|_\infty \leq \eps\}
\end{equation}

CIFAR-10 için yerleşik uzlaşımları izleyerek~\cite{krizhevsky2009learning, croce2021robustbench} birincil pertürbasyon bütçesi olarak $\eps = 8/255 \approx 0{,}031$ kullanıyoruz; epsilon taraması analizi için $\eps \in \{2/255, 4/255, 16/255\}$ değerlerinde ek tam-test PGD-10 değerlendirmeleri yapılmaktadır.

Hem beyaz kutu (tam model erişimi) hem kara kutu (transfer saldırısı) senaryolarında değerlendirme yapıyoruz. Beyaz kutu saldırılarda saldırgan model mimarisi ve parametrelerine tam olarak hâkimdir. Transfer saldırılarında çekişmeli örnekler bir kaynak modelde üretilir ve farklı bir hedef modelde değerlendirilir.

% ----------------------------------------------------------------------------
% 3.2 Model Mimarileri
% ----------------------------------------------------------------------------
\subsection{Model Mimarileri}
\label{subsec:architectures}

\paragraph{CNN Modelleri}

\textbf{ResNet-18}~\cite{he2016deep}: Artık bağlantılı, yaygın kullanılan bir taban CNN; 11,2M parametre içerir. CIFAR-10 için tasarlanmış standart mimariyi (7$\times$7 yerine başlangıçta 3$\times$3 evrişim) kullanıyoruz.

\textbf{WideResNet-28-10}~\cite{zagoruyko2016wide}: ResNet'in 36,5M parametreli daha geniş bir türevi; CIFAR-10'da çekişmeli gürbüzlük için en güncel mimariyi temsil eder. Gowal vd.'nin~\cite{gowal2020uncovering} RobustBench~\cite{croce2021robustbench} üzerinden dağıtılan önceden eğitilmiş gürbüz kontrol noktasını kullanıyoruz; bu model ek etiketsiz veriyle ve bizim standart çekişmeli eğitimimizden daha güçlü bir tarifle eğitildiğinden, eşleşmiş bir karşılaştırma değil bir referans noktası işlevi görmektedir.

\paragraph{Görü Dönüştürücüsü Modelleri}

\textbf{ViT-Tiny}: \texttt{timm} kütüphanesindeki~\cite{rw2019timm} ViT-Tiny mimarisini, girdileri çift doğrusal olarak 224$\times$224'e büyüterek CIFAR-10'a uyarlıyoruz. Mimari 16$\times$16 yamalar (görüntü başına 196 yama), 192 boyutlu gömmeler, her biri 3 dikkat başlı 12 dönüştürücü bloğu ve 4 MLP oranı kullanır; toplam 5,7M parametredir. Model CIFAR-10 üzerinde \emph{sıfırdan} (ImageNet ön eğitimi olmadan) temel veri artırmayla (rastgele kırpma ve yatay çevirme) eğitilmekte, ardından çekişmeli ince ayara tabi tutulmaktadır. Küçük veri kümelerinde ViT gürbüzlüğünün eğitim tarifine duyarlı olduğu bilindiğinden~\cite{debenedetti2023light, mo2022adversarial}, mimari düzeyindeki sonuçlarımız en iyi şekilde eğitilmiş ViT'ler için değil, \emph{bu temel eğitim tarifi altında} geçerli okunmalıdır.

% ----------------------------------------------------------------------------
% 3.3 Saldiri Yontemleri
% ----------------------------------------------------------------------------
\subsection{Saldırı Yöntemleri}
\label{subsec:attack_methods}

\paragraph{FGSM}
Adım boyu pertürbasyon bütçesine eşit tek adımlı saldırı:
\begin{equation}
    x_{adv} = x + \eps \cdot \text{sign}(\nabla_x \mathcal{L}_{CE}(f(x), y))
\end{equation}

\paragraph{PGD}
Adım boyu $\alpha = 2/255$ ve rastgele başlangıçlı yinelemeli saldırı:
\begin{equation}
    x^{t+1} = \Pi_{\mathcal{B}_\eps(x)} \left( x^t + \alpha \cdot \text{sign}(\nabla_{x^t} \mathcal{L}_{CE}(f(x^t), y)) \right)
\end{equation}
burada $x^0 = x + \mathcal{U}(-\eps, \eps)$. Bu çalışma boyunca hem çekişmeli eğitimde hem değerlendirmede PGD-10 (10 yineleme) kullanılmaktadır (Tablo~\ref{tab:main_results}).

\paragraph{AutoAttack}
APGD-CE (çapraz entropi kayıplı Auto-PGD), APGD-T (hedefli DLR kayıplı Auto-PGD), FAB-T (hedefli Hızlı Uyarlanabilir Sınır saldırısı) ve Square (skor tabanlı kara kutu saldırısı) bileşenlerini birleştiren standartlaştırılmış topluluk saldırısı AutoAttack~\cite{croce2020reliable} kullanılmaktadır. Bu topluluk, gürbüzlük değerlendirmesi için standartlaştırılmış güçlü bir kıyaslamadır.

% ----------------------------------------------------------------------------
% 3.4 Savunma Yontemleri
% ----------------------------------------------------------------------------
\subsection{Savunma Yöntemleri}
\label{subsec:defense_methods}

\paragraph{Çekişmeli Eğitim (AT)}
Temiz ve çekişmeli kayıpları birleştiren karma bir çekişmeli eğitim amacı kullanıyoruz:
\begin{equation}
    \mathcal{L} = (1 - \lambda)\,\mathcal{L}_{CE}(f_\theta(x), y) + \lambda\,\mathcal{L}_{CE}(f_\theta(x_{adv}), y)
\end{equation}
burada $\lambda = 0{,}5$, standart sınıflandırma başarımı ile çekişmeli gürbüzlüğü dengeler. İç enbüyükleme, çekişmeli örnekleri PGD-10 ile üretir ($\alpha = 2/255$, $\eps = 8/255$); saldırı üretimi model eğitim kipindeyken yapılır, böylece saldırı üretimi ile parametre güncellemesi tutarlı parti-normalizasyonu istatistikleri görür ve standart çekişmeli eğitim gerçeklemeleri izlenir~\cite{rice2020overfitting}. Pratikte yaygın benimsenen bu karma amaç, gürbüzlük kazanılırken temiz doğruluğun korunmasına yardım eder; yalnızca çekişmeli örneklerle eğitim doğal başarımı düşürebilir. TRADES~\cite{zhang2019theoretically} gibi alternatif formülasyonlar kuramsal temelli ödünleşim mekanizmaları sunsa da, bu çalışma mimari etkileri savunmaya özgü etkileşimlerden yalıtmak için standart çekişmeli eğitime odaklanmaktadır.

% ----------------------------------------------------------------------------
% 3.5 Egitim Protokolu
% ----------------------------------------------------------------------------
\subsection{Eğitim Protokolü}
\label{subsec:training}

\paragraph{Temiz Eğitim}
CNN'ler için SGD eniyileyici; başlangıç öğrenme oranı 0,1, kosinüs tavlama çizelgesi~\cite{loshchilov2016sgdr}, ağırlık sönümü $5 \times 10^{-4}$ ve 200 epok. ViT'ler için AdamW eniyileyici~\cite{loshchilov2017decoupled}; başlangıç öğrenme oranı $10^{-3}$, kosinüs tavlama, ağırlık sönümü 0,05 ve 200 epok.

\paragraph{Çekişmeli Eğitim}
Temiz eğitilmiş kontrol noktalarından başlayarak düşürülmüş öğrenme oranlarıyla ince ayar yapıyoruz: CNN'ler SGD ile $10^{-3}$ öğrenme oranı ve kosinüs tavlamayla en çok 100 epok; ViT'ler AdamW ile aynı öğrenme oranı ve çizelgeyle eğitilir. Gradyanlar küresel $L_2$ normu 10'a kırpılır ve kaybı sonlu olmayan partiler atlanır. Veri artırma, 4 piksel dolgulu rastgele kırpma ve yatay çevirmedir; tüm modeller $[0,1]$ aralığında ham piksel girdisi tüketir (ortalama/std normalizasyonu yoktur) ve saldırılar ile epsilon bütçeleri doğrudan bu piksel uzayında işler. Eğitim parti boyu CNN için 128, ViT için 64'tür (değerlendirme parti boyları analize göre değişir). Erken durdurma 20 sabır ve 0,1 puanlık iyileşme eşiğiyle uygulanır.

\paragraph{Doğrulama Bölmesi ve Tekrarlar}
Model seçimi ve erken durdurma, sabit 2.000 örneklik bir doğrulama bölmesindeki PGD-10 çekişmeli doğruluğuyla yürütülür. Bu bölme bir kez çekilir (tohum 777; indeksler kodla birlikte saklanır), çalışmadaki her model tarafından paylaşılır ve hem temiz hem çekişmeli aşamanın eğitim kümesinden çıkarılır; dolayısıyla hiçbir kontrol noktası, üzerinde eğitildiği veriyle seçilmez. Test kümesine yalnızca nihai art-değerlendirmelerde dokunulur. Tüm boru hattını (temiz ön eğitim + çekişmeli ince ayar) mimari başına üç kez, bağımsız tohumlarla tekrarlıyoruz (ResNet-18 için 1001--1003, ViT-Tiny için 2001--2003) ve koşuları indekse göre eşliyoruz; böylece her mimariler arası karşılaştırma aynı tekrar altında eğitilmiş modeller arasındadır. Tüm tablolar bu üç koşunun ortalamasını ve standart sapmasını raporlar.

Bölüm~\ref{subsec:robustness}'teki bir karşılaştırma için, bu protokol benimsenmeden önce eğitilmiş ve doğrulama bölmesi temiz ön eğitimin parçası olmuş daha eski bir kontrol noktası çiftine de başvuruyoruz. Bunları sonuç olarak değil, açıkça bir karşıtlık olarak raporluyoruz; çünkü iki kontrol noktası kümesi arasındaki fark, tek koşuluk bir sonucun ne kadar hareket edebileceğine dair kendi başına kanıttır.

% ----------------------------------------------------------------------------
% 3.6 Degerlendirme Metrikleri
% ----------------------------------------------------------------------------
\subsection{Değerlendirme Metrikleri}
\label{subsec:metrics}

\paragraph{Temel Metrikler}
Modelleri üç standart metrikle değerlendiriyoruz: temiz doğruluk (pertürbe edilmemiş test görüntülerinde sınıflandırma başarımı), gürbüz doğruluk (çekişmeli örneklerde başarım) ve saldırı başarı oranı (temizde doğru sınıflandırılan görüntülerin saldırı sonrası yanlışa dönen oranı). Beyaz kutu değerlendirmelerde ek olarak \emph{koşullu yanıltma oranını} (temizde doğru sınıflandırılan örneklerin saldırının devirdiği kesri) ve tümleyeni olan koşullu sağkalımı raporluyoruz. Örnek bazındaki kayıtlarımızda hem PGD hem AutoAttack, temizde yanlış sınıflandırılan örnekleri gürbüz olmayan saydığından, mutlak gürbüz doğruluk birebir $\text{gürbüz} = \text{temiz} \times \text{koşullu sağkalım}$ olarak çarpanlarına ayrılır; bu, gürbüzlük farklarını temiz doğruluk bileşeni ile koşullu duyarlılık bileşenine ayırmamızı sağlar. Son olarak \emph{her ikisi doğru} altkümesindeki (karşılaştırılan iki modelin de temizde doğru sınıflandırdığı örnekler) doğrulukları raporluyoruz; bu, ortak bir örnek kümesi üzerinde birebir eşleştirilmiş karşılaştırma sağlar.

\paragraph{Transfer Metrikleri}
Birincil transfer metriğimiz \emph{koşullu yanıltma oranıdır}: hedef modelin temizde doğru sınıflandırdığı örnekler içinde, kaynak modelde üretilen çekişmeli örneklerle yanlışa dönenlerin oranı. Kaynak model $f_S$'ye karşı üretilen $\{x_{adv}^S\}$ örnekleri için hedef $f_T$'ye koşullu transfer oranı:
\begin{equation}
    \text{KYO} = \frac{\sum_i \mathbf{1}[f_T(x_i)=y_i]\,\mathbf{1}[f_T(x_{adv,i}^S) \neq y_i]}{\sum_i \mathbf{1}[f_T(x_i)=y_i]}
\end{equation}

Transfer değerlendirmesini doğru sınıflandırılan örneklerle koşullamak yerleşik pratiktir~\cite{liu2017delving, dong2018boosting, mahmood2021robustness, ravikumar2023trend, zhao2025revisiting} ve mimariler arası karşılaştırmalarda zorunludur: ham yanlış sınıflandırma oranı $\frac{1}{n} \sum_{i} \mathbf{1}[f_T(x_{adv,i}^S) \neq y_i]$ hedefin temiz hatasını içerir; temiz doğruluğu düşük bir hedef, saldırı hiç ek örnek devirmese bile transfere daha ``kırılgan'' görünür. Ham oranı, şeffaflık ve bu karıştırıcının boyutunu nicelemek için ikincil metrik olarak raporluyoruz. Koşullama kümesinin kendisi bir protokol seçimi olduğundan, iki yerleşik alternatifi de raporluyoruz: kaynak ve hedefin ikisinin de doğru sınıflandırdığı örneklerle koşullayan \emph{her ikisi doğru} protokolü~\cite{mahmood2021robustness, ravikumar2023trend} ve saldırının kaynakta beyaz kutu anlamında başarılı olduğu hedef-doğru örneklerle koşullayan \emph{başarılı kaynak} protokolü. Bölüm~\ref{subsec:transfer}'in gösterdiği gibi bu protokoller nitel olarak farklı asimetri sonuçları verebilir; bu yüzden protokolü bir gerçekleme ayrıntısı değil, raporlanan sonucun parçası olarak ele alıyoruz. Transfer oranları tam test kümesinde (10.000 örnek) örnek bazında kayıtla hesaplanır ve örnekler üzerinde yüzdelik bootstrap ile \%95 güven aralıkları verilir (10.000 yeniden örnekleme).

\paragraph{Gradyan Analizi Metrikleri}
Gradyan özelliklerini, büyüklükler değerlendirme parti boyundan bağımsız olsun diye örnek başına kayıplarla hesaplayarak üç tamamlayıcı metrik ailesiyle inceliyoruz. Gradyan normu ($\|\nabla_x \mathcal{L}\|_2$ ve $\|\nabla_x \mathcal{L}\|_\infty$) toplam duyarlılık büyüklüğünü ölçer. Gradyan seyrekliği, ölçekten bağımsız ölçülerle nicelenir --- Hoyer seyrekliği, $|g|$'nin Gini katsayısı ve örnek maksimumunun \%1'inin altındaki bileşen oranı --- çünkü sabit mutlak eşikler gradyan ölçeğini yapıyla karıştırır; Hoyer ve Gini, Hurley ve Rickard'ın~\cite{hurley2009comparing} biçimselleştirdiği ölçek-değişmezliği aksiyomlarını sağlar. Gradyan hizalanması, örnekler arasındaki ortalama ikili mutlak kosinüs benzerliğidir (500 analiz örneğinin tüm çiftleri üzerinden):
\begin{equation}
    \text{Hizalanma} = \frac{1}{N(N-1)} \sum_{i \neq j} \left| \frac{\nabla_{x_i} \mathcal{L} \cdot \nabla_{x_j} \mathcal{L}}{\|\nabla_{x_i} \mathcal{L}\| \|\nabla_{x_j} \mathcal{L}\|} \right|
\end{equation}
Mutlak kosinüsü kullanıyoruz, çünkü paylaşılan duyarlılık \emph{eksenleri} işaretten bağımsız olarak saldırı geometrisi için önemlidir; evrensel pertürbasyon argümanlarının ek olarak gerektirdiği şey işaret tutarlılığı olduğundan işaretli ortalama kosinüsü de raporluyoruz --- bu değer her model için sıfıra yakın çıkmaktadır (Bölüm~\ref{subsec:gradient}). Yüksek mutlak hizalanma, girdiler arasında benzer pertürbasyon eksenlerine işaret eder.

\paragraph{Mekânsal Lokalite Metrikleri}
Seyreklik kaç gradyan bileşeninin büyük olduğunu ölçer ama bunların görüntüde nerede durduğu hakkında bir şey söylemez; bu yüzden mekânsal yapıyı, kanal toplamı enerji haritası $e = \sum_c (\nabla_x \mathcal{L})_c^2$ (toplamı 1'e normalize) üzerinde ayrıca ölçüyoruz. Enerjinin \%50'sini ve \%90'ını toplayan en küçük piksel oranını (küçük değer = daha lokalize), $e$'nin piksel dağılımı olarak Shannon entropisini (düşük değer = daha lokalize) ve enerji bitişik piksellerde kümelendiğinde pozitif olan 4-komşuluklu Moran's $I$'yı raporluyoruz. Dördü de gradyan ölçeğinden bağımsızdır ve örnek başına hesaplanır; böylece iki mimari aynı girdiler üzerinde eşleştirilmiş testlerle karşılaştırılabilir.

\paragraph{Katman Analizi için Jeton Agregasyonu}
Bir dönüştürücü bloğu jeton başına bir vektör üretir; bloğu özetlemek bir seçim gerektirir. Üçünü raporluyoruz: yalnız CLS jetonu, 196 yama jetonunun ortalaması ve tüm jetonların örnek başına düzleştirilmiş birleşimi. CNN'de her artık blok çıktısı mekânsal boyutlar üzerinden düzleştirilir. ViT için üçünü birden raporluyoruz; çünkü Bölüm~\ref{subsec:attention}'ün gösterdiği gibi, farklı kayma büyüklükleri verirler ve minimumu farklı derinliklere yerleştirirler.

\paragraph{Dikkat Metrikleri}
ViT için her bloğun softmax sonrası CLS$\to$yama dikkat satırını başlar üzerinden ortalayarak çıkarıyoruz. Dikkatin ne kadar geniş yayıldığını ölçen Shannon entropisini ve saldırı altındaki yer değiştirmesini (temiz ile çekişmeli dikkat dağılımları arasındaki, 1 ile sınırlı toplam varyasyon uzaklığı) raporluyoruz. İkisi de örnek başına hesaplanıp ortalanır; koşu başına $n=1000$ örnek.

\paragraph{ViT'e Özgü Transfer Saldırısı}
Transfer sonuçlarının mimariden bağımsız saldırılara bağlı olup olmadığını test etmek için kaynak örnekleri ek olarak Jeton Gradyanı Düzenlileştirme (TGR)~\cite{zhang2023tgr} ile üretiyoruz; TGR, geri yayılım sırasında dikkat, QKV ve MLP yollarındaki jeton-gradyan tensörlerinin uç değerli girdilerini sıfırlar ve kalan gradyanı $\gamma=0{,}5$ ile ölçekler. Gerçeklememiz yöntemi CIFAR-10 ViT sarmalayıcımıza ve bu makale boyunca kullanılan bütçeye ($\eps=8/255$, $\alpha=2/255$, 10 adım, momentum 1,0) uyarlamaktadır; özgün ImageNet varsayılanları kullanılmamıştır ve bu yüzden bir uyarlama olarak raporlanmaktadır. Düz MI-FGSM~\cite{dong2018boosting} örnekleri aynı koşuda aynı örnekler üzerinde üretilir; bu, birebir aynı bütçede eşleştirilmiş bir kontrol sağlar.

% ----------------------------------------------------------------------------
% 3.7 Istatistiksel Dogrulama
% ----------------------------------------------------------------------------
\subsection{İstatistiksel Doğrulama}
\label{subsec:statistics}

Üç varyans kaynağını ayırıyor ve her birini geçerli olduğu yerde raporluyoruz. Mimari karşılaştırmalarını ilgilendiren \emph{eğitim koşusu varyansıdır}: her tablo, mimari başına üç bağımsız koşunun ortalamasını ve standart sapmasını verir (Bölüm~\ref{subsec:training}). \emph{Saldırı başlatma varyansı}, PGD-10'un sabit bir altkümede üç saldırı tohumuyla (42, 123, 456) yeniden koşulmasıyla ölçülür; değerlendirilen örnekler ve sıraları sabittir, tohumlar yalnızca saldırının rastgele başlangıcını değiştirir. \emph{Protokol varyansı}, koşullama protokolleri arasındaki yayılımdır ve gürültü değil sonuç olarak ele alınır (Bölüm~\ref{subsec:transfer}). Bölüm~\ref{subsec:statistics_results} büyüklüklerini karşılaştırmaktadır.

Aynı örneklerin iki koşul altında puanlandığı her yerde eşleştirilmiş testler kullanılır: eşleştirilmiş ikili sonuçlar için McNemar'ın kesin testi, üç seyreklik ölçüsü için Holm düzeltmeli Wilcoxon işaretli sıra testi, farklar için eşleştirilmiş bootstrap güven aralıkları (10.000 yeniden örnekleme) ve her-ikisi-doğru transfer farkı için işaret çevirme permütasyon testi (20.000 permütasyon). Eşdeğerlik iddiaları, marjı her kullanımda belirtilen iki tek yönlü test (TOST) ile yapılır. Tüm değerlendirme ve analiz betikleri rastgele tohumları sabitler (tohum 42, CUDA determinizm bayrakları dâhil); dolayısıyla raporlanan her sayı, yayımlanan boru hattından yeniden üretilebilir.

% ============================================================================
% BOLUM 4: DENEYLER
% ============================================================================
\section{Deneyler}
\label{sec:experiments}

% ----------------------------------------------------------------------------
% 4.1 Ana Gurbuzluk Sonuclari
% ----------------------------------------------------------------------------
\subsection{Gürbüzlük Karşılaştırması}
\label{subsec:robustness}

Tablo~\ref{tab:main_results}, değerlendirilen tüm modeller için ana gürbüzlük karşılaştırmasını sunmaktadır. Modellerimize ait her hücre, mimari başına üç bağımsız eğitim koşusunun ortalaması ve bu koşular üzerindeki standart sapmadır; her koşu 10.000 örneklik tam test kümesinde değerlendirilmiştir.

\begin{table}[!t]
\centering
\caption{CIFAR-10'da ana gürbüzlük sonuçları ($\eps = 8/255$); mimari başına üç bağımsız eğitim koşusu üzerinden ortalama$\pm$std. AA, standart takımla AutoAttack anlamına gelmektedir.}
\label{tab:main_results}
\scriptsize
\setlength{\tabcolsep}{2pt}
\begin{tabular}{lcccccc}
\toprule
\textbf{Model} & \textbf{Savunma} & \textbf{Param.} & \textbf{Temiz} & \textbf{FGSM} & \textbf{PGD-10} & \textbf{AA} \\
\midrule
\multicolumn{7}{l}{\textit{CNN Modelleri}} \\
ResNet-18 & Yok & 11,2M & 95,25$\pm$0,15 & 40,29$\pm$1,77 & 0,03$\pm$0,03 & -- \\
ResNet-18 & AT & 11,2M & 85,78$\pm$0,36 & 53,46$\pm$0,08 & 44,11$\pm$0,50 & \textbf{37,93$\pm$0,14} \\
WRN-28-10 & AT$^\dagger$ & 36,5M & 89,48 & 70,91 & 66,92 & 62,76 \\
\midrule
\multicolumn{7}{l}{\textit{Görü Dönüştürücüsü Modelleri}} \\
ViT-Tiny & Yok & 5,7M & 80,09$\pm$0,60 & 6,10$\pm$1,47 & 0,05$\pm$0,05 & -- \\
ViT-Tiny & AT & 5,7M & 73,53$\pm$0,55 & 36,31$\pm$0,42 & 32,69$\pm$0,22 & \textbf{29,14$\pm$0,40} \\
\bottomrule
\multicolumn{7}{l}{\scriptsize $^\dagger$Gowal vd.'nin~\cite{gowal2020uncovering} RobustBench~\cite{croce2021robustbench} üzerinden dağıtılan gürbüz kontrol}\\
\multicolumn{7}{l}{\scriptsize \phantom{$^\dagger$}noktası; tek harici kontrol noktası olduğundan standart sapma yoktur.}\\
\multicolumn{7}{l}{\scriptsize \phantom{$^\dagger$}Temiz/FGSM/PGD-10: yerel tam-test değerlendirmelerimiz; AA: RobustBench}\\
\multicolumn{7}{l}{\scriptsize \phantom{$^\dagger$}raporlu değer (üzerinde AutoAttack yeniden koşulmadı).}
\end{tabular}
\end{table}

CNN her saldırı altında \emph{mutlak} bir gürbüzlük üstünlüğü taşımaktadır: ResNet-18 AT, AutoAttack altında $\%37{,}93\pm0{,}14$ gürbüz doğruluğa ulaşırken ViT-Tiny AT $\%29{,}14\pm0{,}40$'ta kalmaktadır. Bir çiftin iki modeli de aynı 10.000 örnekte, örnek bazında kayıtla değerlendirildiğinden her koşu içinde McNemar'ın eşleştirilmiş testi uygulanabilir; test üç çiftte de CNN lehine kesindir (AutoAttack altında $p < 10^{-84}$, PGD-10 altında $p < 10^{-120}$). Tohum düzeyindeki değişim, çekişmeli eğitilmiş çift için baştan sona küçüktür (raporlanan hiçbir büyüklükte 0,55 puanı aşmaz); dolayısıyla sıralama koşudan koşuya gürültünün ürünü değildir. Savunmasız temel modeller FGSM altında daha çok değişmektedir (1,8 puana kadar), ancak hiçbir iddia onlara dayanmamaktadır. CNN ayrıca daha yüksek temiz doğruluğu korumaktadır ($\%85{,}78$'e karşı $\%73{,}53$; 12,3 puanlık fark); bu, ViT'leri küçük veri kümelerinde sıfırdan eğitmenin güçlüğünü yansıtır; Tablo~\ref{tab:conditional} bu katkıyı koşullu duyarlılıktan ayırmaktadır. AutoAttack gürbüzlüğü iki mimaride de PGD-10'un tutarlı biçimde altındadır (ResNet-18 için 6,2, ViT-Tiny için 3,6 puan); bu, daha güçlü topluluğun değerini doğrular. \%62{,}76'daki WideResNet-28-10 referansı~\cite{gowal2020uncovering}, kapasitenin, ek eğitim verisinin ve daha güçlü tarifin birlikte gürbüzlüğü mimari seçiminin ötesine ne kadar taşıdığını göstermektedir; bu modeli eşleşmiş çiftten görsel ve analitik olarak ayrı tutuyoruz.

\begin{table}[!t]
\centering
\caption{Tam test kümesinde beyaz kutu gürbüzlüğünün koşullu ayrıştırması; üç koşu üzerinden ortalama$\pm$std. Alt blok, aynı büyüklükleri önceki tek koşuluk kontrol noktalarında tekrarlamaktadır.}
\label{tab:conditional}
\scriptsize
\setlength{\tabcolsep}{2.5pt}
\begin{tabular}{lcccc}
\toprule
& \multicolumn{2}{c}{\textbf{Koşullu yanıltma (\%)}} & \multicolumn{2}{c}{\textbf{Her ikisi doğru gürbüz (\%)}} \\
\textbf{Model} & \textbf{PGD-10} & \textbf{AA} & \textbf{PGD-10} & \textbf{AA} \\
\midrule
ResNet-18 AT & \textbf{48,58}$\pm$0,80 & \textbf{55,78}$\pm$0,23 & \textbf{59,86}$\pm$0,35 & \textbf{51,89}$\pm$0,49 \\
ViT-Tiny AT & 55,53$\pm$0,50 & 60,37$\pm$0,33 & 46,11$\pm$0,59 & 41,15$\pm$0,39 \\
\midrule
\multicolumn{5}{l}{\textit{Önceki tek koşu, doğrulama bölmesi temiz ön eğitimde görülmüş ($n$=7.260)}} \\
ResNet-18 AT & 52,15 & 58,16 & \textbf{54,92} & \textbf{48,15} \\
ViT-Tiny AT & 52,33 & \textbf{56,46} & 49,61 & 45,33 \\
\bottomrule
\end{tabular}
\end{table}

Tablo~\ref{tab:conditional}, mutlak skorları temiz doğruluk ve koşullu duyarlılık bileşenlerine ayırmaktadır (Bölüm~\ref{subsec:metrics}). Temizde yanlış sınıflandırılan girdiler gürbüz olmayan sayıldığından ayrıştırma birebirdir: AutoAttack için CNN'de $85{,}78 \times 0{,}442 = 37{,}93$, ViT'te $73{,}53 \times 0{,}396 = 29{,}14$ okunur. İki çarpan da CNN lehinedir: 12,3 puan daha yüksek temiz doğruluktan başlar \emph{ve} çözdüklerinin daha küçük bir kesrini kaybeder (koşullu yanıltma PGD-10 altında $48{,}58$'e karşı $\%55{,}53$; AutoAttack altında $55{,}78$'e karşı $\%60{,}37$). Çiftin iki modelinin de temizde doğru sınıflandırdığı $n=7.061\pm68$ örnekten oluşan ve bu nedenle ortak girdiler üzerinde birebir eşleştirilmiş bir karşılaştırma sağlayan ortak çözülen altkümede CNN 13,8 ve 10,7 puan öndedir (her koşuda McNemar $p<10^{-84}$).

Ayrıştırma, iki çarpan aynı yönü gösterdiğinde bile raporlanmaya değerdir; çünkü aynı yönü göstermek zorunda değillerdir ve aralarındaki pay, ürettikleri sıralamadan daha az kararlıdır. Doğrulama bölmesini temiz ön eğitimin zaten görmüş olduğu ilk kontrol noktalarımız, aynı mutlak sıralamayı ama tersine bir koşullu okuma vermişti: PGD-10 altında ayırt edilemeyen (52,15'e karşı \%52,33) ve AutoAttack altında hafifçe ViT lehine (58,16'ya karşı \%56,46) yanıltma oranları; bu da CNN üstünlüğünün tamamen bir temiz doğruluk etkisi olduğu okumasını destekliyordu (Tablo~\ref{tab:conditional} alt blok). Düzeltilmiş üç koşunun hiçbiri bu okumayı yeniden üretmemekte ve üç koşu birbiriyle, herhangi birinin önceki koşuyla uyuştuğundan çok daha yakından uyuşmaktadır. Tek bir neden yalıtamıyoruz; çünkü düzeltilmiş protokol doğrulama düzenini ve eğitim koşusunu aynı anda değiştirmiştir. Dikkat çekici biçimde basit bir veri miktarı açıklaması geçerli değildir: düzeltilmiş temiz modeller 2.000 daha az görüntü görmelerine karşın öncekilerden \emph{daha iyidir} (ViT için $80{,}09\pm0{,}60$'a karşı $\%77{,}50$). Savunulabilir sonuç, bir mekanizmadan daha dar ve daha yararlıdır: bu türden bir atıf tek eğitim koşusundan çıkarılmamalıdır ve yalnızca mutlak skoru raporlayan bir makale, işi hangi çarpanın yaptığını gizler.

\begin{figure}[!t]
\centering
\includegraphics[width=\columnwidth]{../figures/final_tr/fig1_robustness_comparison.pdf}
\caption{Tam CIFAR-10 test kümesinde $\eps=8/255$'te temiz, PGD-10 ve AutoAttack doğrulukları (üç koşu ortalaması; hata çubukları $\pm$1 std). WRN-28-10'un AutoAttack çubuğu RobustBench tarafından raporlanan değerdir.}
\label{fig:robustness_comparison}
\end{figure}

Şekil~\ref{fig:robustness_comparison} bu gürbüzlük farkını mimariler boyunca görselleştirir. Pertürbasyon büyüklüğünün etkisini anlamak için Şekil~\ref{fig:epsilon_sweep}, $\eps \in \{2/255, 4/255, 8/255, 16/255\}$ bütçelerinde tam test kümesindeki PGD-10 gürbüz doğruluğunu gösterir.

\begin{figure}[!t]
\centering
\includegraphics[width=\columnwidth]{../figures/final_tr/fig2_epsilon_sweep.pdf}
\caption{Pertürbasyon bütçesine ($\eps$) göre PGD-10 gürbüz doğruluğu; tam test kümesi, üç koşu ortalaması ($\pm$1 std bantları). İki AT mimarisi de epsilon arttıkça benzer biçimlerle bozulmakta, CNN tutarlı bir marj korumaktadır.}
\label{fig:epsilon_sweep}
\end{figure}

\begin{figure}[!t]
\centering
\includegraphics[width=\columnwidth]{../figures/final_tr/fig_adversarial_examples.pdf}
\caption{Her modele kendi beyaz kutu PGD-10 saldırısı ve bir CNN$\rightarrow$ViT transfer saldırısı; her biri farklı bir test görüntüsü üzerinde, temiz tahmin, $10\times$ büyütülmüş pertürbasyon ve çekişmeli tahminle. Pertürbasyonlar $\eps$=8/255 ile sınırlıdır ve bu ölçekte görsel olarak silik kalmaktadır.}
\label{fig:adversarial_examples}
\end{figure}

Şekil~\ref{fig:adversarial_examples}, çekişmeli pertürbasyonların ve model tahminleri üzerindeki etkilerinin görsel örneklerini vermektedir.

% ----------------------------------------------------------------------------
% 4.2 Transfer Saldirisi Analizi
% ----------------------------------------------------------------------------
\subsection{Transfer Saldırısı Analizi}
\label{subsec:transfer}

Bu bölüm makalenin merkezi ölçümünü içermektedir. Her tohum çifti için, kaynak modelde PGD-10 çekişmeli örnekleri üretiyor ve bunları tam test kümesinde bir hedefte değerlendiriyoruz. Sonuçlar örnek bazında kaydedildiğinden, aynı saldırı koşusu yeniden hesaplama olmadan birden çok koşullama protokolü altında puanlanabilir: protokoller yalnızca paydaya hangi örneklerin girdiğinde ayrışır (Bölüm~\ref{subsec:metrics}).

\begin{table}[!t]
\centering
\caption{Koşulsuz oran ve üç yerleşik koşullama protokolü altında transfer asimetrisi (PGD-10, $\eps$=8/255, tam test; üç tohum çifti üzerinden ortalama$\pm$std). En sağdaki sütun, önceki tek koşuluk kontrol noktalarında ölçülen farkı tekrarlamaktadır.}
\label{tab:transfer_protocols}
\footnotesize
\setlength{\tabcolsep}{3pt}
\begin{tabular}{lcccc}
\toprule
\textbf{Protokol} & \textbf{CNN$\rightarrow$ViT} & \textbf{ViT$\rightarrow$CNN} & \textbf{Fark} & \textbf{Fark}$^{\ddagger}$ \\
\midrule
Koşulsuz (ham) & 41,02$\pm$0,55 & 27,45$\pm$0,27 & $+$13,57$\pm$0,33 & $+$8,27 \\
Hedef doğru & 19,87$\pm$0,18 & 15,51$\pm$0,59 & $+$4,36$\pm$0,44 & $+$0,63 \\
Her ikisi doğru & 18,25$\pm$0,17 & 9,98$\pm$0,31 & $+$8,27$\pm$0,23 & $+$5,33 \\
Başarılı kaynak & 36,39$\pm$0,76 & 17,02$\pm$0,52 & $+$19,37$\pm$1,27 & $+$11,17 \\
\bottomrule
\multicolumn{5}{l}{\scriptsize $^{\ddagger}$Tek koşu; doğrulama bölmesi temiz ön eğitimde görülmüş.}
\end{tabular}
\end{table}

Paydalar protokoller arasında keskin biçimde değişmektedir. CNN$\rightarrow$ViT ve ViT$\rightarrow$CNN için sırasıyla koşullamasız 10.000/10.000, hedef doğruda 7.353/8.579, her ikisi doğruda 7.061 ortak örnek ve başarılı kaynakta 2.831/3.814'tür. Fark sütununda verilen yayılım, tohum içi eşleştirilmiş farkın standart sapmasıdır; iki yön tohumlar boyunca birlikte değiştiğinden bu değer yön başına yayılımlardan geri hesaplanamaz ve Bölüm~\ref{subsec:statistics_results}'te kullanılan eğitim koşusu referansıdır. Sonuç şudur: ölçülen büyüklüğe model değil protokol hükmetmektedir. Aynı modeller ve aynı veri üzerinde asimetri $+4{,}36\pm0{,}44$ ile $+19{,}37\pm1{,}27$ puan arasında değişmektedir; en büyük ve en küçük protokol tahmini arasındaki tohum-başına ortalama yayılım $15{,}01\pm0{,}84$ puandır (4,4 kat), oysa herhangi bir protokolün içindeki tohum düzeyi standart sapma 1,5 puanın altında kalır. Protokol, istatistiksel hükmü de belirlemektedir. Her iki yönü aynı örnekler üzerinde puanlayan ve bu yüzden eşleştirilmiş çıkarımı destekleyen her ikisi doğru protokolünde fark $8{,}27\pm0{,}24$ puandır; eşleştirilmiş bootstrap GA $[7{,}33; 9{,}22]$, her tohumda işaret çevirme permütasyonu $p<5\times10^{-5}$ ve TOST eşdeğerliği üç marjın üçünde de ($\pm1$, $\pm2$ ve $\pm3$ puan; her birinde $p=1{,}0$) reddedilir. Önceki kontrol noktalarında ise aynı boru hattı, hedef doğru protokolü altında $\pm$2 puanda eşdeğerliği kabul etmişti. Bu sayıların yalnızca birini gören bir okur, mimarilerin farklı olup olmadığı hakkında farklı bir sonuca varırdı.

Bir protokolün adı tek başına bir tanım değildir. Başarılı kaynak protokolü, saldırının kaynakta başarılı olduğu örneklerle koşullamaktadır; bu ifadenin kaynağın temiz girdiyi doğru sınıflandırmasını gerektirip gerektirmediği ise transfer çalışmalarının nadiren yazdığı bir alt seçimdir. Bu makale daha sıkı okumayı raporlamaktadır: beyaz kutu başarısı, kaynağın temiz girdide doğru ve çekişmeli karşılığında yanlış olması demektir. Kaynağın zaten yanlış sınıflandırdığı örnekleri de içeri alan gevşek okumada ölçülen asimetri $+19{,}37$ puandan $+14{,}60$ puana, protokol yayılımı ise $15{,}01$ puandan $10{,}45$ puana inmektedir. Bu $4{,}77$ puanlık kayma, Tablo~\ref{tab:transfer_protocols}'teki her protokolün koşumlar arası standart sapmasının üç katından büyüktür; diğer üç protokol ise tanım gereği değişmemektedir. Bu büyüklükte bir alt seçim protokolün adında görünmediği için, dört koşullama maskesinin tamamı yayımlanan kodda tek bir yerde tanımlanmaktadır.

Kararlı olan yöndür. CNN'de üretilen örnekler ViT'e, dört protokolün dördünde, üç tohumun üçünde ve momentum tabanlı MI-FGSM~\cite{dong2018boosting} altında da ($19{,}07\pm0{,}33$'e karşı $\%15{,}42\pm0{,}52$, hedef doğru) tersinden daha iyi transfer olmaktadır; dolayısıyla etki PGD saldırısının bir artefaktı değildir. Büyüklük, anlamlılığı ve her türlü eşdeğerlik iddiası protokole bağlıdır; işaret değildir.

Yayılımın arkasındaki mekanizma bileşimseldir. Hedefin temizde çözdüğü ama kaynağın çözemediği örnekler, ortak çözülenlerden üç-dört kat daha kırılgandır ($41{,}2\pm0{,}8$ ve $\%59{,}1\pm1{,}7$ yanıltma; ortak kümede $10{,}0$ ve $\%18{,}3$) ve bu marjinal altkümeler, iki model temiz doğrulukta ayrıştığı için boyutça beş kat farklıdır ($1.517\pm102$'ye karşı $291\pm16$). Hedef doğru koşullaması iki marjinal altkümeyi de içeri alır ve asimetrik boyutları iki oranı birbirine doğru çeker; her ikisi doğru koşullaması onları dışarıda bırakır ve yönlü fark yeniden ortaya çıkar. Bu açıklama dört protokolün üçünü kapsar. Başarılı kaynak protokolü farklı bir mekanizmayla sürülür ve iki sürücü ayrıştırıldığında yayılımın neye bağlı olduğu değişir: altı tohum-ve-veri-kümesi kümesinden gelen $18$ yön çiftinde, dört protokollü asimetri yayılımı iki model arasındaki temiz hata farkı büyüdükçe \emph{azalmaktadır} (eğim $-0{,}528$, küme bootstrap GA $[-0{,}664;\ -0{,}418]$); buna karşılık başarılı kaynak protokolü dışarıda bırakılarak hesaplanan yayılım aynı farkla \emph{artmaktadır} (eğim $+0{,}431$, GA $[+0{,}333;\ +0{,}633]$). Başarılı kaynak oranı $18$ çiftin $12$'sinde dördün en küçüğü, $4$ çiftte daha en büyüğüdür; yani $16$ çiftte uçtur Ortalama açıklığı en büyük olan protokol çifti hedef doğru ile başarılı kaynak arasındadır ($23{,}77$ puan) ve bu çift $18$ yönün $13$'ünde en geniştir; kalan beşinde, ki hepsi CIFAR-100'dedir, uçtaki protokol koşulsuz orandır. Kaynak başarısıyla koşullama \emph{kaynağın} kendi gürbüzlüğüne göre seçim yapar; bu, çiftin temiz doğruluk farkının bir fonksiyonu değildir ve dolayısıyla ikinci, kısmen ters yönlü bir terim katmaktadır. Bu uydurmanın havuzu ön kayıtta sabitlenen havuzdur; ön kayıt SVHN çiftini uydurulan noktalara değil bağımsız bir tutarlılık kontrolüne atamaktadır. Temiz hata farkı iki puanın altında kalan ve $x$ ekseninin başka türlü boş olan bir bandına düşen iki SVHN yönü eklendiğinde dört protokollü eğim $-0{,}133$'e gitmekte ve küme bootstrap güven aralığı $[-0{,}570;\ +0{,}469]$ sıfırı içermektedir; üç protokollü eğim ise hem işaretini hem aralığını korumaktadır. Bu duyarlılığı içeri almak yerine raporluyoruz, çünkü dört protokollü eğimin ne kadar ince bir dayanağı olduğunu göstermektedir.

İki kol üç protokolde uyuşmaktadır. Protokol yayılımını temiz hata farkına karşı, farklı biçimlerde hata veren iki yolla kalibre ediyoruz. \emph{Kontrollü} kol, çiftin bir üyesini sabit tutup diğerini kendi eğitim yörüngesi boyunca taramakta; böylece mimari, tohum, veri kümesi ve saldırı bütçesi sabit kalırken fark değişmektedir (sekiz tohum-ve-veri-kümesi kümesinden $116$ nokta). \emph{Gözlemsel} kol ise final modelleri mimariler, tohumlar ve veri kümeleri arasında karşılaştırmakta, yani her şey aynı anda değişmektedir (altı kümeden $18$ yön çifti). Her ikisinde çıkarım yörünge düzeyi küme bootstrap iledir; serbestlik derecesini nokta sayısı değil küme sayısı belirler ve kollar havuzlanmamaktadır. Bu aralıklar için iki uyarı geçerlidir. Altı ile sekiz küme varken yüzdelik küme bootstrap'i kapsamayı düşük tutar; dolayısıyla aşağıdaki aralıklar kesin değil, genişlik için alt sınır olarak okunmalıdır. Ayrıca WideResNet tek ve sabit bir dış kontrol noktasıdır ve altı gözlemsel kümenin altısında da, on sekiz yön çiftinin on ikisinde bulunur; yani bu kümeler tam bağımsız değildir. Onu dışarıda bırakan yalnız-çift duyarlılığı yanında raporlanmaktadır.

Başarılı kaynak protokolü dışarıda bırakıldığında iki kol $+0{,}672$ (GA $[+0{,}601;\ +0{,}726]$) ve $+0{,}431$ (GA $[+0{,}333;\ +0{,}633]$) eğimleri vermektedir: işaret aynı, aralıklar örtüşüyor. Dahil edildiğinde kontrollü kol $+0{,}293$ (GA $[+0{,}206;\ +0{,}489]$), gözlemsel kol $-0{,}567$ (GA $[-0{,}757;\ -0{,}451]$) vermektedir. Bu uyuşmazlık bir kontrol etkisi değil bir bileşim etkisidir: kontrollü kol WideResNet içermez ve gözlemsel kol aynı mimari çiftine kısıtlandığında $+0{,}387$ (GA $[+0{,}075;\ +0{,}576]$) çıkmakta, yani kontrollü kolla hem işaret hem aralık bakımından uyuşmaktadır. Bununla birlikte iki sınır yazılmalıdır: gözlemsel kolun negatif eğimi $x$ ekseninde $7{,}2$ puanlık ölçülmemiş bir boşluğun üzerinden uydurulmuştur, dolayısıyla bir eğilimi izlemek yerine birbirinden ayrık iki kümeyi karşılaştırmaktadır; kontrollü kolun eğimi ise erken kontrol noktaları çıkarıldığında ayakta kalmaktadır (eşiklere göre $+0{,}220$ ile $+0{,}524$ arasında) ve en kararlı nicelik üç protokollü eğimdir ($0{,}659$ ile $0{,}672$).

% ----------------------------------------------------------------------------
% 4.2.1 Uc mimarili kontrol
% ----------------------------------------------------------------------------
\subsubsection{Ham oranın ne kadarı hedefin temiz hatası?}
\label{subsubsec:matrix}

Karıştırıcıyı tartışmak yerine doğrudan nicelemek için analizi, WideResNet-28-10 referansını hem kaynak hem hedef olarak ekleyerek $3\times3$ matrise genişletiyoruz (Tablo~\ref{tab:transfer_matrix}). İki oran arasındaki ilişki bir korelasyondan daha güçlü çıkmaktadır: tek bir olgusal öncüle dayanan bir \emph{özdeşliktir}.

Hedefin temiz hatasına $e$, koşulsuz orana $r_{\text{ham}}$ ve hedef doğru koşullu orana $r_{\text{koş}}$ diyelim. Hedefin temiz girdide zaten yanlış sınıflandırdığı bir örnek pertürbasyondan sonra da yanlış kalıyorsa $r_{\text{ham}} = e + r_{\text{koş}}(1-e)$ olur ve dolayısıyla
\begin{equation}
r_{\text{ham}} - r_{\text{koş}} = e\,(1 - r_{\text{koş}}).
\label{eq:raw_identity}
\end{equation}

Ölçülen kısım öncüldür ve sıkı biçimde tutmaktadır: $36$ köşegen dışı yönde (iki veri kümesi, üç tohum çifti, üç hedef, iki kaynak) temiz girdide yanlış sınıflandırılan bir örneğin saldırı altında da yanlış kalma olasılığı $0{,}989$ ile $1{,}000$ arasındadır. Denklem~\eqref{eq:raw_identity} bunun sonucunda ölçülen her koşulsuz oranı $0{,}41$ puandan küçük bir hatayla yeniden üretmektedir; ortalama mutlak artık $0{,}095$ puandır ve artık, açıklanan sapmanın en çok $\%1{,}39$'unu oluşturmaktadır.

Bu, iddianın konumunu değiştirir ve güçlendirir. Koşulsuz oranın koşullu orandan sapması, verimizde hedefin temiz hatasıyla yalnızca ilişkili değildir; yarım puanın altında bir artık dışında $(1 - r_{\text{koş}})$ ile ölçeklenmiş \emph{o hatanın kendisidir}. Koşulsuz transfer oranları, üç hedefte rastlantısal olarak gözlediğimiz bir düzenlilik olarak değil, \emph{yapısı gereği} ve öncülü sağlayan her çalışmada hedef zayıflığıyla kirlenmektedir. Öncülün kendisi ölçülebilir bir niceliktir ve varsayılmak yerine ölçülmelidir. Temizde yanlış sınıflandırılan bir örneğin saldırı altında yanlış kalma olasılığına $p$ dersek genel biçim $r_{\text{ham}} = e\,(1-p) + r_{\text{koş}}\,(1 - e)$ olur ve $p=1$ olduğunda Denklem~\eqref{eq:raw_identity}'e indirgenir. Çekişmeli eğitilmiş hedeflerde $p$ değerini $0{,}989$ ile $1{,}000$ arasında ölçüyoruz; düzeltme orada puanın onda birkaçı kadar bir hatayla kesindir. Transfer çalışmalarının çoğunun saldırdığı standart eğitilmiş hedeflerde ise kendi önceki koşumlarımızda $p$ değerini $0{,}88$'e kadar inerken gözlemledik; orada özdeşlik bir eşitlik değil bir üst sınırdır. Bu sonucu daha önce bir Pearson korelasyonu olarak ifade etmiştik (üç hedeften gelen altı yön üzerinde $r = 0{,}997$) ve bu ifade sonucu olduğundan zayıf göstermekteydi: bir niceliğin kendi baskın bileşeniyle bire yakın korelasyonu aritmetiktir ve özdeşlik yazıldığında onun gerektirdiği küçük örneklem uyarısına gerek kalmaz.

Veri kümesine özgü kalan şey mekanizma değil orantı çarpanıdır. $(1 - r_{\text{koş}})$ saldırının ne kadar transfer edilebilir olduğuna bağlı olduğundan, uydurulan eğim CIFAR-10'da hedef hatasının puanı başına $0{,}762$, CIFAR-100'de $0{,}656$ puandır. Bu bir doğa sabiti değil, ölçülen hedefler üzerinden bir ortalamadır; özdeşlik başka model havuzlarına uygulanırken varsayılmamalı, yeniden kestirilmelidir. Koşulsuz transfer oranları bu yüzden mimariler arası transfer edilebilirliği ölçmez; transfer edilebilirlik artı hedefin temiz zayıflığını, Denklem~\eqref{eq:raw_identity}'in açıkça gösterdiği bir oranda ölçer.

\begin{table}[!t]
\centering
\caption{$3\times3$ transfer matrisi (PGD-10, $\eps$=8/255, tam test; üç tohum çifti üzerinden ortalama$\pm$std; WideResNet satır ve sütunu tek harici kontrol noktasını kullanır). Ham = koşulsuz yanıltma oranı, Koş. = hedef doğru koşullu oran. Köşegen girdileri beyaz kutu saldırılardır.}
\label{tab:transfer_matrix}
\footnotesize
\setlength{\tabcolsep}{3pt}
\begin{tabular}{llcc}
\toprule
\textbf{Kaynak} & \textbf{Hedef} & \textbf{Ham (\%)} & \textbf{Koş.\ (\%)} \\
\midrule
ResNet-18 AT & ResNet-18 AT & 55,84$\pm$0,53 & 48,52$\pm$0,84 \\
ResNet-18 AT & ViT-Tiny AT & 41,07$\pm$0,58 & 19,93$\pm$0,20 \\
ResNet-18 AT & WRN-28-10 & 21,69$\pm$0,09 & 12,48$\pm$0,10 \\
ViT-Tiny AT & ResNet-18 AT & 27,45$\pm$0,34 & 15,52$\pm$0,64 \\
ViT-Tiny AT & ViT-Tiny AT & 67,34$\pm$0,18 & 55,57$\pm$0,44 \\
ViT-Tiny AT & WRN-28-10 & 18,25$\pm$0,12 & 8,65$\pm$0,14 \\
WRN-28-10 & ResNet-18 AT & 32,63$\pm$0,14 & 21,48$\pm$0,31 \\
WRN-28-10 & ViT-Tiny AT & 39,73$\pm$0,25 & 18,15$\pm$0,28 \\
WRN-28-10 & WRN-28-10 & 33,03 & 25,16 \\
\bottomrule
\end{tabular}
\end{table}

Koşullanmış matris Şekil~\ref{fig:transfer_heatmap}'te verilmektedir. Üçüncü mimari, asimetrinin kendisini de yeniden çerçevelemektedir. İki yabancı kaynak üzerinden ortalandığında gelen koşullu transfer ResNet hedefi için \%18{,}50, ViT hedefi için \%19{,}04 ve WideResNet hedefi için \%10{,}57'dir; bu, her hedefin kendi beyaz kutu koşullu yanıltma oranını ($48{,}52$, $55{,}57$ ve $\%25{,}16$) üç hedef üzerinde $r = 0{,}986$ ile izler. Bu kanıta göre CNN$\rightarrow$ViT üstünlüğü, CNN'de üretilen pertürbasyonların içkin olarak daha transfer edilebilir olduğunun değil, daha zayıf hedefin bir özelliği olarak okunmalıdır: en gürbüz hedef, iki aileden gelen saldırıları neredeyse eşit ölçüde soğurmaktadır. Yalnızca üç hedefle bu korelasyon kesin değil yönlendiricidir ve öyle raporlanmaktadır.

\begin{figure}[!t]
\centering
\includegraphics[width=\columnwidth]{../figures/final_tr/fig3_transfer_heatmap.pdf}
\caption{Koşullu (hedef doğru) transfer yanıltma oranları; üç tohum çifti ortalaması. Bu protokol altında köşegen dışı asimetri $+4{,}4$ puandır; Tablo~\ref{tab:transfer_protocols}, aynı büyüklüğün diğer protokoller altında ne kadar hareket ettiğini gösterir.}
\label{fig:transfer_heatmap}
\end{figure}

% ----------------------------------------------------------------------------
% 4.3 Ikinci veri kumesi (E1)
% ----------------------------------------------------------------------------
\subsection{Protokol Etkisi Genelleniyor mu? İkinci Bir Veri Kümesi}
\label{subsec:cifar100}

Yukarıdaki ölçüm tek bir veri kümesi üzerinde yapılmıştır ve yalnızca CIFAR-10'da var olan bir protokol etkisi genel bir raporlama önerisini desteklemezdi. Bu nedenle tasarımın tamamını CIFAR-100 üzerinde yineledik: aynı mimariler, aynı tehdit modeli, aynı saldırı bütçesi, aynı üç tohumlu yapı, aynı dört protokol. Koşumdan önce üç yönsel ön kestirim ile üç uç nokta bandını salt eklemeli bir belgeye kaydettik; böylece sonuç, ölçüldükten sonra yeniden yorumlanamazdı.

\begin{table}[!t]
\centering
\caption{CIFAR-100 yinelemesi (mimari başına üç bağımsız eğitim koşusunun ortalaması$\pm$std, tam test kümesi, $\eps$=8/255). Her iki model de Bölüm~\ref{subsec:training}'deki protokolle çekişmeli eğitilmiştir. Mutlak gürbüzlük iki mimaride de CIFAR-10'a göre düşüktür ve CNN'in ViT üzerindeki sıralaması değişmemiştir.}
\label{tab:cifar100_results}
\footnotesize
\begin{tabular}{lccc}
\toprule
\textbf{Model} & \textbf{Temiz (\%)} & \textbf{PGD-10 (\%)} & \textbf{AutoAttack (\%)} \\
\midrule
ResNet-18 AT & 63,86$\pm$1,05 & 19,30$\pm$0,32 & 15,04$\pm$0,54 \\
ViT-Tiny AT  & 43,17$\pm$1,10 & 11,15$\pm$0,60 & 8,87$\pm$0,83 \\
\bottomrule
\end{tabular}
\end{table}

Aynı dört protokol boyunca ölçülen CNN$\rightarrow$ViT asimetrisi $+4{,}96\pm1{,}01$ puandan (hedef doğru) $+18{,}53\pm0{,}71$ puana (koşulsuz) uzanmakta; her ikisi doğru protokolünde $+10{,}92\pm0{,}34$, başarılı kaynak protokolünde $+17{,}50\pm0{,}92$ puandır. Tohum başına ortalama yayılım $13{,}83\pm1{,}30$ puandır; CIFAR-10'da bu değer $15{,}01\pm0{,}84$ olduğuna göre daha zor veri kümesi protokol etkisini daraltmamakta, aynı aralıkta bırakmaktadır. Yön, CIFAR-10'da olduğu gibi on iki ölçümün tamamında pozitiftir (dört protokol $\times$ üç tohum çifti). Her ikisi doğru protokolünde eşleştirilmiş fark $10{,}92$ puan, bootstrap güven aralığı $[9{,}47;\ 12{,}35]$, işaret çevirme permütasyon $p$ değeri her tohumda $20{,}000$ permütasyonun çözünürlüğünün altında ve eşdeğerlik sınanan her marjda reddedilmiştir.

Üç ön kestirimin biri doğrulanmış, biri bağlayıcı saydığımız okumada tutmamış, üçüncüsü ise bu tasarımda sınanamamıştır. İşaretin korunacağı ön kestirimi doğrulanmıştır ($12/12$). Protokol yayılımının daha zor veri kümesinde büyüyeceği ön kestirimi, CIFAR-10 yayılımına karşı bir karşılaştırma olarak kaydedilmiş ve o tarihte bu değer $10{,}45$ puan diye anılmıştı. Bu sayısal eşiğe karşı kestirim tutmaktadır ($13{,}83 > 10{,}45$); Bölüm~\ref{subsec:transfer}'in başarılı kaynak tanımı altında CIFAR-10 yayılımının bugünkü değerine karşı ise tutmamaktadır ($13{,}83 < 15{,}01$). Kestirimi kurtaracak okumayı değil, tutmadığını raporluyoruz; çünkü kayıt sayıyı değil CIFAR-10 yayılımını adlandırmıştı. İki veri kümesi protokol etkisini zorluğa göre sıralamak yerine aynı aralığa yerleştirmektedir; yani önerdiğimiz mekanizma etkinin büyüklüğünü öngörmekte, bu iki veri kümesi arasındaki sıralamasını öngörmemektedir. Üçüncüsü, koşulsuz oranın koşullu orandan sapmasının hedefin temiz hatasıyla birlikte büyüyeceğini öngörüyordu. Ölçülen eğim, öngörüldüğü gibi pozitiftir: hedef hatasının puanı başına $+0{,}656$ puan. Ancak CIFAR-100'de üç hedeften ikisinin temiz hatası yalnızca $0{,}21$ puan farklıdır ($36{,}15$ ve $36{,}36$); dolayısıyla korelasyon üç değil \emph{iki} etkin küme üzerinden hesaplanmakta ve değeri (altı yön üzerinde $r = 0{,}931$, hedef düzeyinde $0{,}974$) iki noktalı bir karşıtlığa dayanmaktadır. Ölçüm ön kestirimle \emph{tutarlıdır} ancak onu \emph{sınamaz}; bu yüzden üçüncü bir doğrulama olarak değil, böyle raporlanmaktadır.

Kayıtlı uç noktalardan biri tutmamıştır. Çekişmeli eğitilmiş ViT için temiz doğruluğu $\%33$ ile $40$ bandında öngörmüş, $\%43{,}17\pm1{,}10$ ölçtük. Band, çekişmeli eğitimin temiz doğruluk bedelinin iki mimaride aynı oranda ölçekleneceği varsayımıyla türetilmişti. Ölçekleşmiyor: ViT çekişmeli eğitim altında temiz ön eğitim doğruluğunun $0{,}885$'ini korurken ResNet $0{,}802$'sini korumuştur. Protokol iddiasının dayandığı iki çekişmeli uç nokta ise bantlarının içinde kalmıştır (PGD-10 öngörü $\%10$ ile $14$, ölçüm $11{,}15$; AutoAttack öngörü $\%8$ ile $11$, ölçüm $8{,}87$).

Sınıf bileşimi asimetriyi üretmemektedir. Koşullama, paydaya hangi örneklerin gireceğini değiştirdiği için sınıf karışımını da değiştirir; ilkece bir asimetri bileşim artefaktı olabilir. Her yönün oranını bir bileşim bileşenine ve bir oran bileşenine ayrıştırıyoruz. Hedef doğru protokolünde bileşim etkisi üç tohum çiftinde $-1{,}249$, $-1{,}330$ ve $-1{,}065$ puandır; yani asimetrinin $\%12{,}9$ ile $18{,}6$'sı kadar ve \emph{negatiftir}: sınıf karışımı dengelenseydi asimetri küçülmez, büyürdü. Başarılı kaynak protokolünde etki küçüktür ve üç tohumun üçünde de pozitiftir ($+0{,}527$, $+0{,}435$ ve $+0{,}382$ puan; asimetrinin \%$2{,}1$ ile $3{,}1$'i kadar). Koşulsuz ve her ikisi doğru protokollerinde tam olarak sıfırdır; bu, ayrıştırmanın geçmesi gereken aritmetik sağlama testidir.

Bu sonuçla birlikte okunması gereken iki niteleme vardır. Birincisi, iki çekişmeli band en yakın tohum düzeyi kenarında $0{,}22$ puanlık (AutoAttack) ve $0{,}48$ puanlık (PGD-10) marjlarla tutmuştur. Bu marjlar, \emph{aynı veri kümesindeki} kontrol noktası seçimi gürültüsünden küçüktür: eğitim koşusu sabit tutulup CIFAR-100 yörüngelerinde erken durdurma sabrı ile eğri yumuşatması tarandığında raporlanan PGD-10 doğruluğu $0{,}18$ ile $1{,}83$ puan arasında oynamaktadır (ortalama $0{,}86$). Aynı dokuz hücreli tarama CIFAR-10'da doğrulama bölmesi başına $0{,}54$ ile $2{,}83$ puan vermekte; Bölüm~\ref{subsec:statistics_results}'te anılan on sekiz hücreli değer ise daha da büyüktür, çünkü bölmeyi de değiştirmektedir ve CIFAR-100 protokolünde o boyut yoktur. Bantlar tutmuştur, ama kaydedildikleri veri kümesinin seçim gürültüsüne eşit ya da ondan küçük marjlarla; okur ön kaydı buna göre tartmalıdır. İkincisi, çekişmeli eğitilmiş kontrol noktasını seçmekte kullanılan doğrulama bölmesi, temiz ön eğitimde de kullanılmış olan bölmedir. Bu tercih, iki veri kümesi karşılaştırılabilir kalsın diye CIFAR-10 protokolünden bilinçli olarak devralınmıştır ve Bölüm~\ref{subsec:statistics_results} bunun neye mal olup neye mal olmadığını raporlamaktadır.

% ----------------------------------------------------------------------------
% 4.4 Ucuncu veri kumesi: SVHN, mimarilerin esitlendigi rejim
% ----------------------------------------------------------------------------
\subsection{Mimarilerin Eşitlendiği Üçüncü Bir Veri Kümesi}
\label{subsec:svhn}

CIFAR-10 ve CIFAR-100, iki mimariyi temiz doğrulukta birbirinden hayli uzağa yerleştirmektedir (yaklaşık $11$ ve $21$ puan). SVHN~\cite{netzer2011reading} öyle değildir (karşılaştırma için CIFAR-100 mutlak değerleri Tablo~\ref{tab:cifar100_results}'te verilmektedir): aynı tarifle çekişmeli eğitildiklerinde ResNet $\%94{,}27\pm0{,}70$, ViT $\%92{,}42\pm0{,}56$ temiz doğruluğa ulaşmakta, yani fark $1{,}85$ puanda kalmaktadır. PGD-10 altındaki gürbüz doğruluklar $55{,}52\pm0{,}64$ ve $52{,}30\pm0{,}10$'dur. Tezimizin en keskin öngörüde bulunduğu rejim budur ve öngörü, önceki bölümleri okuyanın bekleyeceği yönde değildir.

\begin{table}[!t]
\centering
\caption{SVHN'de dört protokol altında transfer asimetrisi (PGD-10, $\eps$=8/255, tam test kümesi $n$=26.032; iki tohum çiftinin ortalaması$\pm$std; yalnız iki mimari olduğu için üçüncü mimari kontrolü yoktur). CIFAR-10 ve CIFAR-100'den farklı olarak işaret dört satırda aynı değildir.}
\label{tab:svhn_protocols}
\footnotesize
\setlength{\tabcolsep}{3pt}
\begin{tabular}{lccc}
\toprule
\textbf{Protokol} & \textbf{CNN$\rightarrow$ViT} & \textbf{ViT$\rightarrow$CNN} & \textbf{Fark} \\
\midrule
Koşulsuz (ham) & 31,89 & 31,49 & $+$0,39$\pm$0,08 \\
Hedef doğru & 26,32 & 27,33 & $-$1,00$\pm$0,22 \\
Her ikisi doğru & 25,36 & 25,70 & $-$0,35$\pm$0,01 \\
Başarılı kaynak & 61,38 & 58,69 & $+$2,70$\pm$0,40 \\
\bottomrule
\end{tabular}
\end{table}

Mimariler eşitlendiğinde işarete protokol karar vermektedir. Dört protokolün ikisi CNN üstünlüğü, ikisi ViT üstünlüğü bildirmektedir (Tablo~\ref{tab:svhn_protocols}). Bu ayrım keyfî değildir. Yönlü iddialar için benimsediğimiz iki protokol, hedef doğru ile her ikisi doğru, birbiriyle uyuşmakta ve ikisi de ViT lehine çıkmaktadır; işaret değişimi bunları koşulsuz orandan ve başarılı kaynak oranından ayırmaktadır. Üstelik koşulsuz orandaki çevrilme Denklem~\eqref{eq:raw_identity}'in aritmetik sonucudur: burada ViT hedefinin temiz hatası CNN'inkini yaklaşık iki puan aşmaktadır ve bu, $-1{,}00$'lik koşullu farkı $+0{,}39$'luk ham farka taşımaya yetmektedir. Dolayısıyla ham oran raporlamak yayımlanmış bir karşılaştırmanın işaretini çevirebilir ve özdeşlik bunun ne zaman olacağını önceden söylemektedir. İki yönü özdeş örnekler üzerinde puanladığı için birincil saydığımız her ikisi doğru protokolünde fark $-0{,}35$ puan, eşleştirilmiş bootstrap güven aralığı $[-0{,}77;\ +0{,}06]$, işaret çevirme permütasyon $p$ değeri $0{,}105$ ve TOST eşdeğerliği, denediğimiz en dar marj dâhil ($\pm1$ puan, $p=0{,}0007$) üç marjın üçünde de \emph{kabul edilmektedir}. Yalnızca başarılı kaynak oranını raporlayan bir çalışma SVHN'de $2{,}70$ puanlık bir CNN üstünlüğü duyururdu; yalnızca hedef doğru oranını raporlayan bir çalışma ise $1{,}00$ puanlık bir ViT üstünlüğü. İkisi de aritmetik olarak doğru olurdu ve hiçbiri diğerince yeniden üretilemezdi.

Bu davranış mekanizmaya aykırı değil, onunla tutarlıdır. SVHN'de protokol yayılımı $3{,}70\pm0{,}62$ puandır; CIFAR-10'da $15{,}01\pm0{,}84$, CIFAR-100'de $13{,}83\pm1{,}30$ puandır; temiz doğruluk farkları ise sırasıyla $1{,}85$, $\approx11$ ve $\approx21$ puandır. İlişki üçü arasında tekdüze değildir: en büyük fark en büyük yayılımı vermemekte, CIFAR-100 CIFAR-10'un bir miktar altında kalmaktadır. Veri kümelerini burada ayıran şey, eşitlenmeye yakın çift ile diğerleri arasındaki uzaklıktır; eşitlenmemiş ikisi arasındaki sıralama değil. Yayılım yeterince küçüldüğünde ise asimetrinin kendisini aşamamakta ve işaret, mimarilerin değil protokolün bir özelliği hâline gelmektedir. Bölüm~\ref{subsec:transfer}'de kararlı diye raporladığımız yön, \emph{iki modelin temiz doğrulukta ayrıştığı rejimde} kararlıdır ve her iki CIFAR veri kümesi de o rejimde bulunmaktadır.

SVHN, önceden sabitlenmiş indirgenmiş bir yapılandırmayla koşulmuştur: üç yerine iki tohum çifti, $50$ epok çekişmeli eğitim, AutoAttack yok ve üçüncü mimari yok. Sonuncusu, Bölüm~\ref{subsubsec:matrix}'teki $3\times3$ karıştırıcı kontrolünün burada kurulamayacağı anlamına gelir ve onun SVHN içi bir sürümünü kurmuyoruz.

% ----------------------------------------------------------------------------
% 4.5 Tehdit modelinin degistirilmesi
% ----------------------------------------------------------------------------
\subsection{Protokol Etkisi Tehdit Modeli Değişince Ayakta Kalıyor mu?}
\label{subsec:l2}

Buraya kadarki her ölçüm tek bir tehdit modelinin içinde yapılmıştır. Protokol etkisinin $\Linf$'ye özgü bir artefakt olup olmadığını sınamak için aynı CIFAR-10 kontrol noktalarını $\eps=0{,}5$'lik bir $L_2$ bütçesi altında yeniden değerlendirdik ve üç yönsel ön kestirimi, herhangi bir $L_2$ çıktısına bakmadan önce kaydettik. Modeller $\Linf$ ile eğitilmiştir; dolayısıyla bu ölçüm $L_2$ gürbüzlüğünü değil \emph{ölçüm eksenini} yoklamakta ve aşağıdaki mutlak değerler $L_2$ ile eğitilmiş referanslarla karşılaştırılamamaktadır.

\begin{table}[!t]
\centering
\caption{Aynı kontrol noktaları, $L_2$ bütçesi ($\eps$=0,5; üç tohum çiftinin ortalaması$\pm$std). PGD-L2 tam test kümesini, AutoAttack-L2 ise ön kayıtlı $5.000$ örneklik bir altkümeyi kullanır; bu yüzden temiz sütunu tam küme değerinden bir miktar farklıdır ve ikisi aynı örneklerden gelmez.}
\label{tab:l2_results}
\footnotesize
\begin{tabular}{lcccc}
\toprule
\textbf{Model} & \textbf{Temiz} & \textbf{PGD-L2} & \textbf{Temiz (5k)} & \textbf{AA-L2} \\
\midrule
ResNet-18 AT & 85,78 & 63,78$\pm$0,38 & 85,33 & 60,61$\pm$0,42 \\
ViT-Tiny AT & 73,53 & 51,72$\pm$0,23 & 73,61 & 49,85$\pm$0,20 \\
\bottomrule
\end{tabular}
\end{table}

Kayıtlı üç ön kestirimin üçü de tutmaktadır. Protokol yayılımı $L_2$ altında $10{,}92$ puandır; $\Linf$ altında $15{,}01$ idi. Yani norm değişince küçülmekte, ama ön kayıtın koyduğu iki puanlık tabanın çok üstünde ve koşumlar arası standart sapmadan bir mertebe büyük kalmaktadır. İşaret on iki ölçümün tamamında pozitiftir ve koşulsuz oranın koşullu orandan sapması hâlâ hedefin temiz hatasıyla artmaktadır (eğim $+0{,}914$; korelasyon yerine eğimi raporluyoruz, çünkü on sekiz yön yalnızca yedi ayrık hedef hatası değeri almakta ve bunların altısı aynı sabit WideResNet'ten gelmektedir, dolayısıyla bunlar üzerinde bir korelasyon Bölüm~\ref{subsec:transfer}'in anlattığı biçimde kesinliği abartırdı). Dört protokolün (Tablo~\ref{tab:l2_results}) bildirdiği asimetriler $+13{,}03\pm0{,}81$ (koşulsuz), $+2{,}12\pm0{,}18$ (hedef doğru), $+4{,}07\pm0{,}36$ (her ikisi doğru) ve $+12{,}06\pm0{,}23$ (başarılı kaynak) şeklindedir. Mutlak gürbüzlük her iki saldırı altında da CNN'in ViT üzerindeki sıralamasını korumakta, McNemar $p<10^{-57}$ her tohum çiftinde sağlanmaktadır.

Eğim, özdeşliği örneklem dışında da doğrulamaktadır. Denklem~\eqref{eq:raw_identity}, eğimin hedefler üzerinden ortalanmış $1-r_{\text{koş}}$ değerine eşit olmasını öngörür. Koşullu oranlar bu $L_2$ bütçesi altında $\Linf$'dekinden çok daha düşüktür (kabaca $\%2$ ile $9$'a karşı $\%9$ ile $22$), dolayısıyla özdeşlik \emph{daha dik} bir eğim öngörür; ölçülen $0{,}914$ tam da olması gereken yerdedir; aynı modellerde $\Linf$ altında bu değer $0{,}762$ idi. Orantı çarpanı bir sabit değil, saldırı bütçesinin bir özelliğidir; Bölüm~\ref{subsubsec:matrix} bunu zaten belirtmektedir.

% ----------------------------------------------------------------------------
% 4.6 Gradyan Analizi
% ----------------------------------------------------------------------------
\subsection{Gradyan Karakteristikleri}
\label{subsec:gradient}

Girdi duyarlılığındaki mimari farkları betimlemek için, iki mimarinin gradyan özelliklerini, ana ViT hattımızdaki $32\rightarrow224$ büyütmenin etkisini yalıtan CIFAR'a özgü bir ViT kontrolüyle ($32\times32$ girdi, $4\times4$ yama, büyütme yok; aynı çekişmeli amaç işleviyle, ancak Bölüm~\ref{subsec:training}'te anlatılan sabit bölme protokolü benimsenmeden önce eğitilmiş) birlikte analiz ediyoruz. Tablo~\ref{tab:gradient}, ölçekten bağımsız seyreklik ölçüleriyle (Bölüm~\ref{subsec:metrics}) temel istatistikleri özetler; gradyanlar, normlar değerlendirme parti boyuna bağlı olmasın diye örnek başına kayıplarla hesaplanır. Oradaki hizalanma, tüm $500\times499/2$ örnek çifti üzerindeki ortalama mutlak kosinüs benzerliğidir ve ölçekten bağımsız seyreklik metriklerindeki her ResNet ile ViT farkı, her koşuda örnek başına eşleştirilmiş Wilcoxon testlerinde anlamlıdır; Holm düzeltmeli $p < 10^{-11}$ ve eşleştirilmiş Cohen $d$ değeri 0,32 ile 1,23 arasındadır.

\begin{table}[!t]
\centering
\caption{Örnek başına kayıp gradyanlarından hesaplanan gradyan karakteristikleri (koşu başına $n$=500 test örneği; üç koşu üzerinden ortalama$\pm$std).}
\label{tab:gradient}
\footnotesize
\setlength{\tabcolsep}{2.5pt}
\begin{tabular}{lccc}
\toprule
\textbf{Metrik} & \textbf{ResNet-18 AT} & \textbf{ViT-Tiny AT} & \textbf{ViT (yerel)$^\ddagger$} \\
\midrule
Hoyer seyrekliği & \textbf{0,4928}$\pm$0,0120 & 0,4561$\pm$0,0055 & 0,405 \\
Gini katsayısı & \textbf{0,6498}$\pm$0,0089 & 0,6099$\pm$0,0051 & 0,571 \\
Maks.\ \%1 altı kesir & \textbf{\%34,8}$\pm$2,1 & \%26,8$\pm$0,9 & \%17,9 \\
Hizalanma (mutlak) & 0,0378$\pm$0,0009 & \textbf{0,0562}$\pm$0,0004 & 0,079 \\
Hizalanma (işaretli) & 0,0014 & 0,0005 & -- \\
\bottomrule
\multicolumn{4}{l}{\scriptsize $^\ddagger$CIFAR'a özgü ViT kontrolü ($4\times4$ yama, büyütme yok; 5,4M parametre);}\\
\multicolumn{4}{l}{\scriptsize \phantom{$^\ddagger$}sabit bölme protokolünden önceye ait, test kümesinde seçilmiş tek kontrol noktası;}\\
\multicolumn{4}{l}{\scriptsize \phantom{$^\ddagger$}yalnızca büyütme etkisinin işareti için eklenmiştir, nicel karşılaştırma için değil.}
\end{tabular}
\end{table}

Üç ölçekten bağımsız seyreklik ölçüsünün üçünde de CNN gradyanları ViT gradyanlarından daha seyrektir ve sıralama, tohum düzeyi standart sapmaları farktan en az üç kat küçük kalarak her koşuda yeniden üretilmektedir. Aynı 500 örnek iki modelde de ölçüldüğünden farklar örnek başına eşleştirilmiş istatistiklerle test edilir: üçü de üç koşuda da anlamlıdır (Wilcoxon işaretli sıra, baştan sona Holm düzeltmeli $p<10^{-11}$; eşleştirilmiş Cohen $d$ 0,32 ile 1,23 arasında). Sıralama mimarinin değil çekişmeli eğitimin ürünüdür: aynı koşuların temiz eğitilmiş karşılıklarında daha seyrek model \emph{ViT'tir} (Hoyer $0{,}380\pm0{,}009$'a karşı $0{,}347\pm0{,}002$); yani çekişmeli eğitim CNN gradyanlarını çok daha güçlü seyrekleştirmektedir~\cite{chalasani2020concise} ile tutarlı olarak. Yerel çözünürlüklü ViT kontrolü her ölçüde büyütülmüş ViT'ten daha az seyrektir; dolayısıyla sıralama $32\rightarrow224$ büyütmenin artefaktı değildir. Örnekler arası hizalanma ViT'te mutlak kosinüste daha yüksektir ve temiz eğitilmiş çift bu farkı zaten göstermektedir ($0{,}058$'e karşı $0{,}029$); bu da onu savunmanın değil mimarinin bir özelliği yapar. \emph{İşaretli} ortalama kosinüs ise iki modelde de sıfıra yakındır ($\leq0{,}0014$): örneklerin paylaştığı şey duyarlılık eksenleridir, tutarlı yönler değil; dolayısıyla hizalanma farkı ucuz evrensel pertürbasyonlar için kanıt sayılamaz.

Bu gradyan sonuçları, analizimizin hangi kontrol noktası kümesinin kullanıldığına duyarsız olan tek parçasıdır: Tablo~\ref{tab:gradient}'teki her büyüklük, önceki tek koşuluk kontrol noktalarında ölçülen değeri raporlanan standart sapmalar içinde yeniden üretmektedir; Bölüm~\ref{subsec:robustness}'in koşullu ayrıştırmasının aksine.

% ----------------------------------------------------------------------------
% 4.3.1 Seyreklik mekansal lokalite degildir
% ----------------------------------------------------------------------------
\subsubsection{Seyreklik mekânsal lokalite değildir}
\label{subsubsec:spatial}

Seyreklik ölçüleri kaç gradyan bileşeninin büyük olduğunu yanıtlar; bu bileşenlerin görüntüde bir arada durup durmadığını değil. Ayrım önemlidir; çünkü daha-seyrek-CNN sonucu mekânsal bir okumaya davet çıkarır; gradyanların nesne bölgelerinde daha ``lokalize'' ya da ``yoğunlaşmış'' olduğu düşünülür. Bu okumayı belirli bir çalışmaya atfetmiyoruz; sınadık, çünkü bir seyreklik sonucunun doğal yorumu budur ve kendi önceki şekillerimiz de onu teşvik etmişti. Bu okuma, bu ölçülerin destekleyemeyeceği bir iddia içerir. Bu yüzden mekânsal yapıyı aynı gradyanlar üzerinde, kanal toplamı enerji haritasının dört ölçekten bağımsız betimleyicisiyle doğrudan ölçüyoruz: gradyan enerjisinin \%50'sini ve \%90'ını taşıyan en küçük piksel oranı, normalize enerji haritasının mekânsal entropisi ve enerji bitişik piksellerde kümelendiğinde pozitif olan 4-komşuluklu Moran's $I$.

\begin{table}[!t]
\centering
\caption{Girdi gradyanlarının mekânsal lokalitesi (koşu başına $n$=500 örnek; üç koşu üzerinden ortalama$\pm$std). Eşleştirilmiş farklar aynı örnekler üzerinde hesaplanır. Çıkarım koşum düzeyindedir: $p$, üç koşumun ortalama farkları üzerinde tek örneklem $t$ testidir ve dört ölçüye Holm düzeltmesi uygulanmıştır; böylece çıkarımın birimi iddianın birimiyle eşleşir. Düşük alan ve entropi daha lokalize, yüksek Moran's $I$ daha kümeli demektir.}
\label{tab:spatial}
\footnotesize
\setlength{\tabcolsep}{2.5pt}
\begin{tabular}{lcccc}
\toprule
\textbf{Ölçüt} & \textbf{ResNet-18} & \textbf{ViT-Tiny} & \textbf{Fark} & \textbf{$p$} \\
\midrule
\%50 enerji alanı & 0,0378$\pm$0,0017 & 0,0397$\pm$0,0010 & $-$0,0019 & 0,54 \\
\%90 enerji alanı & 0,2345$\pm$0,0078 & 0,2570$\pm$0,0067 & $-$0,0224 & 0,40 \\
Mekânsal entropi (nat) & 5,2014$\pm$0,0387 & 5,2524$\pm$0,0197 & $-$0,0509 & 0,54 \\
Moran's $I$ & 0,4141$\pm$0,0132 & 0,3941$\pm$0,0024 & $+$0,0200 & 0,47 \\
\bottomrule
\end{tabular}
\end{table}

Sonuç negatiftir (Tablo~\ref{tab:spatial}). Dört ölçütün hiçbiri koşum düzeyinde anlamlı fark göstermez. En büyükleri olan, enerjinin \%90'ını taşıyan alan, CNN için görüntü alanının iki puanı kadar \emph{daha küçük} bir alana işaret eder ve düzeltmeden sonra ayakta kalmaz; bu da $p<10^{-11}$'de anlamlı olan Hoyer farkının yanında marjinal bir etkidir. Tek bir koşum içindeki örnek bazlı testler bu ölçülerde çok daha küçük $p$ değerlerine ulaşır, ama farklı bir soruyu yanıtlarlar, yani iki sabit kontrol noktasının bu 500 görüntüde farklı olup olmadığını; üstelik dört ölçünün ikisinde işaret koşumlar boyunca kararlı değildir. Komşuluk kümelenmesi iki mimaride esasen özdeştir. İki ölçü ailesi bu yüzden uyuşmamaktadır ve uyuşmazlık bilgilendiricidir: çekişmeli eğitilmiş CNN gradyanları kütlelerinin daha büyük kısmını daha az piksele yerleştirir, ama bu pikseller ViT'inkilerden mekânsal olarak daha kümeli değildir. CNN gradyanlarının mekânsal olarak lokalize olduğu betimlemeleri ölçümlerimizce desteklenmemektedir; kendi iddialarımızı seyreklikle sınırlıyoruz.

\begin{figure}[!t]
\centering
\includegraphics[width=\columnwidth]{../figures/final_tr/fig4b_gradient_distribution.pdf}
\caption{İki mimaride örnek başına gradyan $L_2$ normlarının dağılımı. Aynı girdiler için enerji haritaları Şekil~\ref{fig:gradient_heatmap}'te verilmektedir.}
\label{fig:gradient_analysis}
\end{figure}

Şekil~\ref{fig:gradient_analysis}, bu gradyan farkları için görsel kanıt sağlar.

\begin{figure}[!t]
\centering
\includegraphics[width=\columnwidth]{../figures/final_tr/fig4_gradient_comparison.pdf}
\caption{Aynı girdiler üzerinde CNN ve ViT gradyan enerji haritalarının karşılaştırması. Farklar Tablo~\ref{tab:gradient}'teki seyreklik ölçüleriyle nicelenir; Tablo~\ref{tab:spatial}, görünür yapının mekânsal kümelenme farkına karşılık gelmediğini gösterir.}
\label{fig:gradient_heatmap}
\end{figure}

% ----------------------------------------------------------------------------
% 4.4 Oznitelik Kaymasi
% ----------------------------------------------------------------------------
\subsection{Görü Dönüştürücülerinde Öznitelik Kayması}
\label{subsec:attention}

Şimdi çekişmeli hasarın derinlik boyunca nerede biriktiğini ölçüyor, temiz eğitilmiş ViT'lerin katman bazlı temsil analizlerini~\cite{raghu2021vision} çekişmeli eğitilmiş ortama taşıyoruz. Bu ölçümde genellikle örtük bırakılan iki tasarım seçimi önemli çıkmaktadır, bu yüzden ikisini de değiştiriyoruz: örneklem boyutu (tek küçük parti yerine koşu başına $n=1000$) ve bir dönüştürücü bloğunu özetlemekte kullanılan jeton agregasyonu. ViT için yalnız CLS jetonunu, 196 yama jetonunun ortalamasını ve tüm jetonların düzleştirilmiş birleşimini; CNN için mekânsal boyutlar üzerinden düzleştirilmiş her artık blok çıktısını raporluyoruz.

\begin{table}[!t]
\centering
\caption{Her modelin kendi beyaz kutu PGD-10 saldırısı altında blok bazlı temizden çekişmeliye kosinüs benzerliği (koşu başına $n$=1000; üç koşu üzerinden ortalama$\pm$std). ViT sütunları yalnızca blok jetonlarının nasıl toplandığında ayrışır.}
\label{tab:feature_degradation}
\footnotesize
\setlength{\tabcolsep}{2.5pt}
\begin{tabular}{lccc|lc}
\toprule
\multicolumn{4}{c|}{\textbf{ViT-Tiny AT}} & \multicolumn{2}{c}{\textbf{ResNet-18 AT}} \\
\textbf{Blok} & \textbf{CLS} & \textbf{Yama ort.} & \textbf{Tüm jetonlar} & \textbf{Blok} & \textbf{Kosinüs} \\
\midrule
1 & 0,9964 & 0,9990 & 0,9975 & layer1.0 & 0,9951 \\
2 & 0,9893 & 0,9954 & 0,9911 & layer1.1 & 0,9918 \\
3 & 0,9785 & 0,9928 & 0,9821 & layer2.0 & 0,9879 \\
4 & 0,9709 & 0,9877 & 0,9734 & layer2.1 & 0,9853 \\
5 & 0,9652 & 0,9870 & 0,9719 & layer3.0 & 0,9753 \\
6 & 0,9596 & 0,9873 & 0,9712 & layer3.1 & 0,9607 \\
7 & 0,9497 & 0,9870 & 0,9698 & layer4.0 & \textbf{0,8777} \\
8 & 0,9414 & \textbf{0,9840} & \textbf{0,9663} & layer4.1 & 0,9108 \\
9 & 0,9353 & 0,9850 & 0,9675 & & \\
10 & \textbf{0,9343} & 0,9877 & 0,9696 & & \\
11 & 0,9352 & 0,9912 & 0,9753 & & \\
12 & 0,9779 & 0,9955 & 0,9884 & & \\
\bottomrule
\end{tabular}
\end{table}

Bundan iki sonuç çıkmaktadır. Birincisi, iki mimari hasarın \emph{birikip birikmediğinde} değil \emph{nerede} biriktiğinde ayrışmaktadır. ViT'te kayma kademeli ve sığdır: benzerlik 8 ile 10. bloklar civarındaki bir minimuma doğru düzenli biçimde azalır ve çıkışa doğru toparlanır; öznitelik normları baştan sona temiz değerlerinin \%0,7'si içinde kalır, yani pertürbasyon temsili yeniden ölçeklemek yerine döndürmektedir. CNN'de hasar yoğunlaşmıştır: ilk üç aşama az kaybeder, ardından son artık aşamanın ilk bloğu $-\%13{,}14\pm1{,}89$'luk norm daralmasıyla birlikte $0{,}8777\pm0{,}0181$'e çöker; son blok temsili kısmen onarır ($0{,}9108\pm0{,}0060$, norm $+\%5{,}84\pm3{,}53$). CNN'de çekişmeli hasar geç, yerel ve büyüklük değiştiren bir olaydır; ViT'te dağıtık, yalnızca yön değiştiren bir dönmedir. Önceki kontrol noktalarımızda CNN profilinin son bloğa dek monoton düştüğünü, toparlanma göstermediğini not ediyoruz; dolayısıyla ViT profilinin aksine CNN eğrisinin biçimi tarife bağlı okunmalıdır.

İkincisi, agregasyon seçimi ölçülen profili önemli ölçüde değiştirmektedir. Kendi minimumlarında CLS yolu, yama jetonu ortalamasının yaklaşık dört katı kadar kaymakta ($0{,}9343$'e karşı $0{,}9840$; $1-\cos$ cinsinden oran $4{,}1$) ve iki özet, minimumu farklı derinliklere bile yerleştirmektedir (blok 10'a karşı blok 8). Sınıflandırıcının okuduğu şey CLS jetonu olduğundan bu bir ayrıntı değil bir sonuçtur: pertürbasyon sınıflandırma yoluna yoğunlaşırken ortalama yama temsili görece korunmaktadır. Bu ayrıca, farklı makalelerden katman bazlı kayma eğrilerinin, agregasyon belirtilmedikçe doğrudan karşılaştırılamayacağı anlamına gelir.

Blok bazlı profiller Tablo~\ref{tab:feature_degradation}'te verilmektedir. Üçüncü bir ölçüm negatiftir. Çekişmeli pertürbasyon, CLS dikkat dağılımının entropisini ölçülebilir biçimde değiştirmemektedir: koşu başına $n=1000$ örnek üzerinde, katman başına ortalama entropi değişimi mutlak değerce en çok $0{,}005$ nattır (katman ortalamaları 4,9 nat civarındayken) ve tohum düzeyi standart sapma aynı mertebededir. Bunu reddedememe olarak değil eşdeğerlik olarak ifade ediyoruz: katman ortalamasının yüzde biri olan ve veriden değil ölçünün ölçeğinden sabitlenen $\pm 0{,}05$ nat marjıyla TOST, on iki katmanın on ikisinde de eşdeğerliği kabul etmektedir; en zayıfı $p=0{,}006$. Derinlikle büyüyen şey dikkat haritasının \emph{yer değiştirmesidir}: ilk katmanda $0{,}028$ toplam varyasyondan son katmanda $0{,}084$'e çıkmaktadır; toplam varyasyon 1 ile sınırlı olduğundan mutlak olarak küçüktür. Saldırı dikkatı hafifçe yeniden yönlendirmekte, ama ne genişletmekte ne keskinleştirmektedir. Bunu açıkça işaretliyoruz; çünkü bu analizin 8 örneklik bir figür partisiyle hesaplanan önceki bir sürümü, daha büyük ölçümün desteklemediği görünür bir dikkat etkisi izlenimi vermişti; çekişmeli saldırının dikkati bozduğuna dair her iddia, dile getirilmeden önce bu örneklem boyutunda sınanmalıdır. Bunu literatürdeki belirli bir çalışmaya karşı değil, kendi ön sürümümüze karşı söylüyoruz.

Son olarak kayma büyüklüğü saldırı başarısını öngörmemektedir. Örnekler, saldırının tahmini devirip devirmediğine göre bölündüğünde, ViT'te devrilen ve devrilmeyen örneklerin son blok kosinüs benzerliği istatistiksel olarak ayırt edilemezdir ($0{,}9895$'e karşı $0{,}9877$); CNN'de beklenen yönde küçük bir fark vardır ($0{,}9058$'e karşı $0{,}9143$). Öznitelik kayması, bu yüzden, temsilin nasıl hareket ettiğinin bir betimlemesidir; tahminin değişip değişmeyeceğinin yeterli istatistiği değildir.

\begin{figure}[!t]
\centering
\includegraphics[width=\columnwidth]{../figures/final_tr/fig5_attention_comparison.pdf}
\caption{Çekişmeli saldırı altında ViT katmanları boyunca CLS jetonu dikkat haritaları (softmax sonrası, başlar üzerinden ortalanmış). Her satır farklı bir katmanı (1, 6, 12) gösterir; sütunlar temiz dikkat, çekişmeli dikkat ve farklarıdır. Gösterilen tek bir örnektir; dikkat kaymasındaki derinlik eğilimi metinde koşum başına $n=1000$ örnek üzerinden nicel olarak verilmektedir.}
\label{fig:attention_layers}
\end{figure}

Nicel ölçümleri tamamlayıcı olarak Şekil~\ref{fig:attention_layers}, temsilî bir örnek için 1., 6.\ ve 12.\ katmanlardaki CLS dikkat haritalarını görselleştirir. Görselleştirme yukarıdaki yer değiştirme ölçümüyle tutarlıdır, çünkü geç katmanlar erken katmanlardan daha çok hareket eder; ancak bu türden tek örnekli dikkat panellerinin kolayca fazla yorumlandığına dikkat çekiyoruz: 1.000 örnek üzerindeki toplam entropi değişimi sıfırdır, yani görünür farklar yayılmayı değil yeniden yönlendirmeyi yansıtır. Yama pertürbasyonları altındaki dikkat kayması etkileri~\cite{gu2022vision} ve dikkat mekanizmasını doğrudan hedefleyen saldırılar~\cite{fu2022patchfool}, burada incelenen $L_\infty$ sınırlı pertürbasyonlardan farklı bir rejimde işler.

\begin{figure}[!t]
\centering
\includegraphics[width=\columnwidth]{../figures/final_tr/fig5b_attention_entropy.pdf}
\caption{Temiz ve çekişmeli girdiler için CLS dikkat dağılımının katman başına entropisi; koşum başına $n=1000$ örnek, üç koşum ortalaması ve $\pm$1 std bandı. Dikkat haritalarının kendisi Şekil~\ref{fig:attention_layers}'te verilmektedir.}
\label{fig:feature_degradation}
\end{figure}

Şekil~\ref{fig:feature_degradation}, temsilî bir örneğin temiz ve çekişmeli CLS dikkatini katman başına entropi eğrileriyle birlikte gösterir; iki entropi eğrisi üst üste binmektedir; yukarıda raporlanan sıfır sonucun görsel biçimi budur.

\begin{figure}[!t]
\centering
\includegraphics[width=\columnwidth]{../figures/final_tr/fig_tsne_features.pdf}
\caption{Temiz örnekler (daireler) ve her modelin kendi beyaz kutu PGD-10 çekişmeli örnekleri (üçgenler) için sondan bir önceki katman özniteliklerinin gösterim amaçlı t-SNE gömmesi; ResNet-18 (sol) ve ViT-Tiny (sağ). t-SNE yalnızca yönlendirme içindir; nicel karşılaştırma metindedir.}
\label{fig:tsne_features}
\end{figure}

Şekil~\ref{fig:tsne_features}, sondan bir önceki katman öznitelik uzaylarının gösterim amaçlı bir t-SNE gömmesini verir. t-SNE gömmeleri hiperparametrelere duyarlı olduğundan ve nitel aşırı yorumu davet ettiğinden, tüm iddiaları özgün öznitelik uzayında hesaplanan metriklere dayandırıyoruz (koşu başına $n$=500; her model beyaz kutu saldırıya uğratılmış; üç koşu üzerinden ortalama$\pm$std). Saldırı, sınıf yapısını en güçlü biçimde \emph{CNN'de} çökertmektedir: temizde eğitilmiş 5-NN etiket tutarlılığı $\%87{,}3\pm2{,}6$'dan $\%51{,}3\pm1{,}8$'e düşer, siluet skoru $0{,}196\pm0{,}010$'dan $-0{,}032\pm0{,}004$'e iner, sınıf içi/sınıflar arası uzaklık oranı $0{,}511\pm0{,}013$'ten $0{,}847\pm0{,}003$'e çıkar ve göreli sınıf merkezi kayması $0{,}238\pm0{,}012$'dir. ViT'in sondan bir önceki uzayı temiz girdilerde bile zayıf sınıf ayrımı gösterir (siluet $-0{,}021\pm0{,}005$; 5-NN tutarlılığı $\%61{,}1\pm1{,}4$, saldırıda $\%48{,}5\pm1{,}4$'e düşer) ve merkezleri çok daha az kayar ($0{,}083\pm0{,}004$). Ölçülebilir ifade şudur: CNN'in iyi ayrışmış temiz öznitelik uzayı beyaz kutu saldırıyla ağır biçimde bozulurken, ViT'in sınıf yapısı baştan daha zayıftır ve bu toplu ölçülerde daha az hareket eder; bu, gürbüzlük hakkında değil geometri hakkında bir ifadedir, zira daha yüksek gürbüz doğruluğu yine de CNN korumaktadır.

% ----------------------------------------------------------------------------
% 4.5 ViT'e ozgu transfer saldirisi
% ----------------------------------------------------------------------------
\subsection{ViT'e Özgü Bir Transfer Saldırısı Tabloyu Değiştiriyor mu?}
\label{subsec:tgr}

Buraya kadarki tüm transfer sonuçları, ViT'i genel türevlenebilir bir model olarak ele alan gradyan saldırılarını kullanmaktadır. Doğal bir itiraz, dönüştürücü yapısı için tasarlanmış saldırıların daha iyi transfer ettiği bildirildiğinden, ViT yönünün bu seçim tarafından cezalandırıldığıdır. İtirazı, kaynak modelde geri yayılım sırasında dikkat, QKV ve MLP yollarındaki uç değerli jeton gradyanlarını bastıran Jeton Gradyanı Düzenlileştirme (TGR)~\cite{zhang2023tgr} ile doğrudan test ediyoruz. Gerçeklememiz bir uyarlamadır: özgün yöntem standart eğitilmiş ImageNet ViT'lerini hedefler; biz aynı mekanizmayı CIFAR-10 ViT-Tiny sarmalayıcımıza uyguluyor ve saldırı bütçesini ImageNet varsayılanları yerine makalenin geri kalanıyla eşliyoruz ($\eps=8/255$, $\alpha=2/255$, 10 adım, momentum 1,0). Kontrol olarak aynı koşuda düz MI-FGSM örnekleri üretilir; böylece karşılaştırma özdeş örnekler ve özdeş bütçe üzerinde eşleştirilmiştir.

\begin{table}[!t]
\centering
\caption{Eşleşmiş bütçede TGR~\cite{zhang2023tgr} ve düz MI-FGSM altında ViT$\rightarrow$CNN transferi (koşu başına $n$=10.000; üç koşu üzerinden ortalama$\pm$std).}
\label{tab:tgr}
\footnotesize
\begin{tabular}{lcc}
\toprule
\textbf{Ölçüt (\%)} & \textbf{TGR} & \textbf{MI-FGSM} \\
\midrule
Kaynakta beyaz kutu (ham) & 51,56$\pm$0,80 & \textbf{66,43}$\pm$0,29 \\
Transfer, hedef doğru & 12,70$\pm$0,41 & \textbf{15,71}$\pm$0,60 \\
Transfer, her ikisi doğru & 7,93$\pm$0,42 & \textbf{10,08}$\pm$0,30 \\
\bottomrule
\end{tabular}
\end{table}

TGR bu rejimde transferi iyileştirmemekte; kötüleştirmektedir (Tablo~\ref{tab:tgr}). Düzenlileştirme kaynak tarafında tasarlandığı gibi davranır ve 14,9 puanlık beyaz kutu gücünden vazgeçer; ancak bu takas hedefte karşılığını bulmaz: koşullu transfer hedef doğru protokolünde 3,0, her ikisi doğruda 2,2 puan düşer. Her-ikisi-doğru altkümesindeki eşleştirilmiş McNemar testleri her koşuda MI-FGSM lehine kesindir ($p = 6{,}1\times10^{-10}$, $1{,}1\times10^{-27}$, $7{,}0\times10^{-34}$); yani bu bir tohum etkisi değildir.

Sonucu dar tutuyoruz. Kurulumumuz, TGR'nin tasarlandığı kurulumdan her biri önemli olabilecek üç noktada ayrılır: hedefler standart değil çekişmeli eğitilmiştir, veri yerel ImageNet değil model içinde büyütülen $32\times32$ CIFAR-10 görüntüleridir ve kaynak 5,7M parametreli bir ViT-Tiny'dir. Destekleyebildiğimiz iddia, TGR'nin transfer üstünlüğünün bu rejime taşınmadığıdır; yöntemin genel olarak başarısız olduğu değil. Bu makale açısından önemli sonuç şudur: Bölüm~\ref{subsec:transfer}'te raporlanan CNN$\rightarrow$ViT asimetrisi, mimariden bağımsız saldırılar kullanmanın bir artefaktı değildir; ViT yönüne dönüştürücüye özgü bir saldırı vermek onu güçlendirmemekte, zayıflatmaktadır.

% ----------------------------------------------------------------------------
% 4.6 Varyans Kaynaklari
% ----------------------------------------------------------------------------
\subsection{Varyans Kaynakları}
\label{subsec:statistics_results}

Yukarıdaki karşılaştırmaları dört varyans kaynağı ilgilendirmektedir. Saldırı başlatma stokastisitesi ihmal edilebilirdir: PGD-10'u sabit 1.000 örneklik altkümede üç saldırı tohumuyla yeniden koşmak temiz doğruluğu hiç değiştirmez (yükleyici karıştırmaz) ve gürbüz doğruluğu CNN'de $0{,}05$, ViT'te $0{,}00$ puanlık standart sapmayla oynatır. Eğitim koşusu varyansı daha büyüktür ama yine küçüktür: mimari başına üç bağımsız koşuda, çekişmeli eğitilmiş çift için standart sapma temiz doğrulukta en çok $0{,}55$, PGD-10'da $0{,}50$ puandır (Tablo~\ref{tab:main_results}).

Protokol etkisini doğrudan bu iki sayıyla karşılaştırmak bir ölçek hatası olurdu; çünkü protokol etkisi transfer asimetrisi üzerinde, yukarıdaki iki sayı ise mutlak doğruluk üzerinde ölçülmektedir. Bu nedenle protokol ve eğitim koşusu etkilerini \emph{aynı} nicelik olan köşegen dışı asimetri üzerinde ve mutlak birimlerle niceliyoruz. Sabit bir tohum çiftinde yalnızca koşullama protokolünü değiştirmek asimetriyi $15{,}01\pm0{,}84$ puan oynatmaktadır. Bu değer tohum-başına açıklıkların ortalamasıdır ve burada protokol ortalamalarının açıklığıyla çakışmaktadır, çünkü uç protokoller üç tohumun üçünde de aynıdır: üstte başarılı kaynak, altta hedef doğru. CIFAR-100'de ise iki değer ayrışmaktadır ($13{,}83$'e karşı $13{,}58$), çünkü orada uç protokol tohumlar arasında değişmektedir. Aynı asimetriyi, her iki mimariyi sıfırdan yeniden eğitmek ise yalnızca $0{,}23$ ile $1{,}48$ puanlık standart sapmayla oynatmaktadır; değer hangi protokol altında ölçüldüğüne bağlıdır (Tablo~\ref{tab:transfer_protocols}, Fark sütunu). Makalenin dayandığı sıralama budur: on beş puanlık bir ölçüm etkisine karşı $1{,}5$ puanın altında bir eğitim etkisi.

Bu iki yayılımın oranını bilerek manşete koymuyoruz. Üç koşumla payda iki serbestlik derecesi taşır; dolayısıyla herhangi bir tek oran için $\chi^2$ aralığı çok geniştir: dördü içinde en büyük koşumlar arası standart sapmaya sahip başarılı-kaynak protokolünde oranın $\%95$ aralığı $0{,}8$ ile $9{,}9$ arasında uzanır ve \emph{birimi içerir}. Tahmin olarak değil duyarlılık olarak raporlandığında oran, yayılım ölçüleri ve referans protokoller boyunca $5{,}2$ ile $32{,}6$ arasında değişmekte ve birincil saydığımız her-ikisi-doğru protokolünde benzeri-benzeriyle standart sapma karşılaştırmasında $28{,}3$ olmaktadır. İki uyarı oranı daha da sınırlar. En büyük asimetriyi veren protokol aynı zamanda en büyük koşumlar arası standart sapmayı taşımaktadır; yani pay ile payda bağımsız değildir. Ayrıca üç koşumun tamamı, eğitimden önce çekilen tek doğrulama bölmesini paylaştığından bu standart sapma bölmeler arası varyansı dışarıda bırakır; bu da oranı bir üst sınır yapar. Bu aralığın genişliğini ve iki uyarıyı da atlatan şey mutlak karşılaştırmadır; iddiayı taşıyan da kat değeri değil odur. Bu sıralama makalenin savının nicel hâlidir: protokol belirtilmemişse okur, gerçek bir mimari farkı bir ölçüm seçiminden ayırt edemez.

Dördüncü kaynak, genellikle örtük bırakılan çevrimdışı kontrol noktası seçimi protokolüdür. Bir eğitim yörüngesi tamamen sabit tutulup yalnızca son kontrol noktasının nasıl seçildiği değiştirildiğinde, yani hangi doğrulama bölmesi, hangi erken durdurma sabrı ve argmaks öncesinde doğrulama eğrisinin yumuşatılıp yumuşatılmadığı (on sekiz bileşim), raporlanan PGD-10 doğruluğu ResNet'te tohum başına $2{,}62$ ile $2{,}85$, ViT'te $1{,}58$ ile $2{,}09$ puan oynamaktadır. Bu hücrelerde diskteki ağırlıklar özdeştir; yalnızca içlerinden birini seçen kural farklıdır. Bu yayılımın yaklaşık yarısı üç boyuttan yalnız birinden gelmektedir. Argmaks öncesinde doğrulama eğrisini yumuşatmak, üç ResNet koşumunun üçünde ve üç ViT koşumunun ikisinde test PGD-10 doğruluğunu monoton biçimde düşürmektedir; yani yansız bir seçenek değil, sistematik olarak daha kötü bir kuraldır. Çalışmaların fiilen değiştirdiği bölme ve sabır seçimleriyle sınırlandığında yayılım ResNet'te tohum düzeyi standart sapmanın $1{,}09$ katına, ViT'te $1{,}42$ katına inmektedir. İki değeri de veriyoruz: büyüğü belgelenmemiş bir ön işleme adımının ekleyebileceğini, küçüğü belgelenmiş bir protokolün yine de açık bıraktığını ölçmektedir. Bu yayılımı mutlak birimlerle raporluyor ve bilinçli olarak tohum düzeyi yayılıma oran biçiminde vermiyoruz; çünkü o oran, on sekiz hücreden hangisinin referans alındığına bağlıdır: on sekizinin tamamı tarandığında ResNet'te $0{,}67$ ile $3{,}38$, ViT'te $0{,}63$ ile $7{,}38$ arasında değişmekte, bazı referanslarda birin altına inip karşılaştırmayı tersine çevirmektedir.

Hemen yanına bir karşı ağırlık konmalıdır. CIFAR-100'de kendi protokolümüzün seçtiği epok üç ResNet tohumu arasında $25$ epok kaymakta, buna karşılık ortaya çıkan test PGD-10 standart sapması $0{,}32$ puanda kalmaktadır. Seçim \emph{yolu} oynaktı, seçilen \emph{sonuç} değildi. İki gözlem birbirinden ayrı genellenemez: seçim tercihleri raporlanan bir sayıyı eğitim tohumundan daha fazla oynatabilir, neredeyse hiç oynatmayabilir de; hangisinin gerçekleşeceği mimarinin değil, yörüngenin platosunun bir özelliğidir. Raporlama açısından pratik sonuç iki durumda da aynıdır: seçim kuralı sonucun bir parçasıdır ve yöntem bölümüne yazılmalıdır.

\textbf{Hesaplama Ayrıntıları:} Burada raporlanan eğitim koşuları ve analizler NVIDIA RTX 5090 (32GB) üzerinde yürütülmüştür; Tablo~\ref{tab:conditional}'deki karşılaştırmada kullanılan önceki tek koşuluk kontrol noktaları RTX 5060 Ti (16GB) üzerinde eğitilmişti. Tüm koşular PyTorch 2.6.0, CUDA 12.8 ve deterministik cuDNN ayarlarıyla yapılmıştır. AutoAttack değerlendirmesi, standart yapılandırmayla (APGD-CE, APGD-T, FAB-T, Square) 10.000 örneklik tam test kümesinde, örnek bazında kayıtla, tohumlanmış 1.000 örneklik parçalar hâlinde koşulmuştur. Üç tohumlu eğitim boru hattı toplam 21,7 GPU-saat gerektirmiştir.

% ============================================================================
% BOLUM 5: TARTISMA
% ============================================================================
\section{Tartışma}
\label{sec:discussion}

% ----------------------------------------------------------------------------
% 5.1 Ana Bulgular
% ----------------------------------------------------------------------------
\subsection{Ölçümlerin Gösterdikleri}
\label{subsec:findings}

\paragraph{Protokol varyansı, karşılaştırılan büyüklüklere hükmediyor}

Merkezî sonuç bir yayılım sıralamasıdır. Transfer asimetrisinin kendisi üzerinde ölçüldüğünde, her iki mimariyi sıfırdan yeniden eğitmek onu $0{,}23$--$1{,}48$ puanlık standart sapmayla oynatırken, yalnızca koşullama protokolünü değiştirmek $10{,}45\pm0{,}76$ puan oynatmaktadır (Bölüm~\ref{subsec:statistics_results}). İki yayılımı kat değeri olarak değil mutlak birimlerle veriyoruz; çünkü üç koşumla o oranın aralığı bir manşeti taşıyamayacak kadar geniştir. Farklı bir nicelik üzerinde ölçek fikri vermek üzere: bir saldırıyı yeni bir rastgele başlangıçla yeniden koşmak mutlak gürbüz doğruluğu yaklaşık $0{,}05$ puan, yeniden eğitmek ise çekişmeli eğitilmiş çifti en çok $0{,}55$ puan oynatır. Ölçüm seçiminin on puan oynattığı, yeniden eğitmenin ise bir puan oynattığı bir büyüklük mimarilerin kararlı bir özelliği değildir ve protokol olmadan raporlandığında neredeyse hiçbir şey iletmez. Bu, koşullamaya karşı bir argüman değildir --- doğru sınıflandırılan örneklerle koşullamak yerleşik pratiktir~\cite{liu2017delving, dong2018boosting, ravikumar2023trend} ve gerekli olmaya devam etmektedir --- seçimin dipnota değil sonuca ait olduğuna dair bir argümandır.

Dört protokol birbirinin yerine geçmez ve her birinin savunulabilir bir kullanımı vardır. Koşulsuz oran mimari karşılaştırması için hiç kullanılmamalıdır: $3\times3$ matrisimiz, koşullu orandan sapmasının hedefin temiz hatasının $(1 - r_{\text{koş}})$ ile ölçeklenmiş hâline cebirsel olarak eşit olduğunu gösterir; yani transfer edilebilirliği değil, bilinen bir oranda hedef zayıflığını ölçer. Hedef doğru protokolü bu karıştırıcıyı giderir ama kaynağın kendisinin çözemediği örnekleri kabul eder; bu marjinal altkümeler aşırı kırılgan ve asimetrik boyutlu olduğundan iki yönü, temiz doğruluk farkına bağlı bir miktarda birbirine doğru çeker. Her ikisi doğru protokolü iki yönü özdeş örnekler üzerinde puanlar ve dördü içinde eşleştirilmiş çıkarımı destekleyen tek protokoldür; birincil tahmin olarak onu bu yüzden alıyoruz. Başarılı kaynak protokolü farklı bir soruyu yanıtlar --- zaten işe yaramış bir saldırı ne kadar taşınır --- ve daha büyük değeri, kaynak başarısıyla koşullamanın güçlü pertürbasyonları seçmesini yansıtır.

\paragraph{Yön yalnızca modeller eşitlenmemişken kararlıdır}

Dört protokol, üç eğitim koşusu ve iki saldırı ailesi (PGD-10 ve MI-FGSM) boyunca, CNN'de üretilen çekişmeli örnekler \emph{her iki CIFAR veri kümesinde} ViT'e tersinden daha iyi transfer olmaktadır. Aynı tarifle çekişmeli eğitimin iki mimariyi $11$ ya da $21$ puan değil yalnızca $1{,}85$ puan ayırdığı SVHN'de ise bu artık doğru değildir: iki protokol CNN üstünlüğü, iki protokol ViT üstünlüğü bildirmekte ve birincil her ikisi doğru protokolünde fark $-0{,}35$ puan olup eşdeğerlik kabul edilmektedir (Bölüm~\ref{subsec:svhn}). Yön, mimarilerin değil eşitlenmemiş rejimin bir özelliğidir. Değişen, etkinin ne kadar büyük göründüğü ve bir eşdeğerlik testinin onu kabul edip etmediğidir. Üç hedefli analizimiz bu yön için CNN pertürbasyonlarının özel olmasını gerektirmeyen bir okuma önermektedir: gelen transfer her hedefin kendi beyaz kutu kırılganlığını izler (üç hedef üzerinde $r = 0{,}986$) ve en gürbüz hedef iki aileden gelen saldırıları neredeyse eşit soğurur. Bu okumaya göre asimetri, hangi modelin daha zayıf hedef olduğuna dair bir ifadedir; dönüştürücüye özgü bir saldırının~\cite{zhang2023tgr} ViT yönünü kurtaramaması da bununla tutarlıdır. Üç hedefle korelasyon yönlendiricidir; kesinleştirmek daha büyük bir model havuzu gerektirir.

\paragraph{Manşet sayının içindeki atıf, manşetin kendisinden daha kararsız}

Koşullu ayrıştırma, mutlak gürbüzlüğü temiz doğruluk çarpanı ile koşullu duyarlılık çarpanına ayırır. Düzeltilmiş koşularımızda ikisi de CNN lehinedir; ayrıştırma ile manşet uyuşur. Önceki kontrol noktalarımızda uyuşmuyorlardı: koşullu duyarlılık PGD-10 altında eşleşik, AutoAttack altında hafifçe ViT lehine görünüyordu. Mutlak sıralama iki kontrol noktası kümesinde de özdeş olduğundan, yalnızca mutlak doğruluğu raporlayan bir makale hiçbir değişiklik görmezdi; tek koşudan ayrıştırma raporlayan bir makale ise tekrarlanmayan mekanizma düzeyinde bir sonuç çıkarırdı. Pratik ders şudur: mekanizma iddiaları sıralama iddialarından daha çok eğitim koşusu ister ve ikisine aynı kanıt ağırlığı verilmemelidir.

\paragraph{Ölçümden sağ çıkamayan üç iddia}

Seyreklik mekânsal lokalite değildir. Çekişmeli eğitilmiş CNN gradyanları kütlenin daha büyük kısmını daha az bileşene yoğunlaştırır (Hoyer $0{,}4928$'e karşı $0{,}4561$, eşleştirilmiş $p<10^{-11}$); ama aynı gradyanların dört mekânsal betimleyicisi, bu bileşenlerin görüntüdeki dizilişinde karşılık gelen bir fark göstermez. CNN gradyanlarının mekânsal olarak lokalize olduğu betimlemeleri bu yüzden ayrı kanıt ister; seyreklik ölçüleri~\cite{hurley2009comparing} bunu sağlamaz.

Çekişmeli pertürbasyon CLS dikkat entropisini ölçülebilir biçimde değiştirmez. $n=1000$'de katman başına entropi değişimi en çok $0{,}005$ nattır; dikkat haritaları ise küçük ve derinlikle artan bir miktarda yer değiştirir. Yama saldırılarıyla elde edilen dikkat kayması sonuçları~\cite{gu2022vision, fu2022patchfool} farklı bir tehdit modeline aittir ve ölçüm olmadan $L_\infty$ sınırlı pertürbasyonlara genellenmemelidir.

ViT'e özgü bir transfer saldırısı burada ViT kaynaklı transferi iyileştirmez. TGR~\cite{zhang2023tgr} 14,9 puan beyaz kutu gücüne mal olmakta ve bunun hiçbirini hedefte geri vermemektedir. Bunu rejime özgü bir negatif sonuç olarak raporluyoruz --- hedeflerimiz çekişmeli eğitilmiş, verimiz CIFAR ölçeğinde --- ve tam da bu türden tek bir test edilmemiş varsayımın, protokole duyarlı bir karşılaştırmaya taşıyacağı yük yüzünden.

\paragraph{Ayakta kalanlar}

İki davranışsal fark bağımsız kontrol noktası kümelerinde tekrarlanmakta ve protokol sorusundan etkilenmemektedir. Gradyan yapısı farklıdır: CNN girdi gradyanları üç ölçekten bağımsız ölçünün üçünde de daha seyrektir, bu sıralamayı mimari değil çekişmeli eğitim yaratır (temiz eğitilmiş çiftte tersidir) ve örnekler arası hizalanma ViT'te mutlak kosinüste daha yüksekken işaretli ortalama iki modelde de sıfıra yakındır --- örnekler yönleri değil duyarlılık eksenlerini paylaşır; ölçümümüzü bir evrensel pertürbasyon iddiasından~\cite{moosavi2017universal} ayıran ayrım budur. Derinlik profilleri farklıdır: CNN'de çekişmeli hasar son artık aşamaya sıkışmış geç, yerel, büyüklük değiştiren bir olay; ViT'te dağıtık, yalnızca yön değiştiren bir dönmedir. ViT'in ölçülen profilinin jeton agregasyonuna bağlı olması --- CLS yolunun yama ortalamasının üç katı kayması --- katman bazlı çalışmalar için kendi başına bir raporlama gerekliliğidir.

% ----------------------------------------------------------------------------
% 5.2 Oneriler
% ----------------------------------------------------------------------------
\subsection{Raporlama Gereklilikleri}
\label{subsec:implications}

Yukarıdaki ölçümler, mimariler arası gürbüzlük çalışmaları için kısa ve somut bir liste doğurmaktadır.

\begin{enumerate}
\item \textbf{Koşullama protokolünü belirtin ve eşleştirilmiş olanları yeğleyin.} Paydaya hangi örneklerin girdiğini raporlayın. Yönlü bir iddia yapılacaksa, eşleştirilmiş çıkarım mümkün olsun diye iki yönü özdeş örnekler üzerinde puanlayan bir protokol kullanın.
\item \textbf{Koşulsuz transfer oranlarını mimari karşılaştırması olarak raporlamayın.} Hedefler temiz doğrulukta ayrışıyorsa ham oran, transfer edilebilirlik ile hedef zayıflığının \emph{bilinen} bir karışımıdır: $r_{\text{ham}} - r_{\text{koş}} = e\,(1 - r_{\text{koş}})$. Bu uydurulmuş bir ilişki değil bir özdeşlik olduğundan, düzeltme yayımlanmış ham oranlara hiçbir şey yeniden koşulmadan uygulanabilir.
\item \textbf{Yalnızca skoru değil, ayrıştırmayı raporlayın.} Mutlak gürbüzlük, temiz doğruluk ile koşullu duyarlılığın çarpımına birebir ayrışır; manşeti hangi çarpanın taşıdığı işin ilginç kısmıdır ve skordan geri çıkarılamaz.
\item \textbf{Sıralama iddialarını mekanizma iddialarından ayırın.} Sıralamalarımız kontrol noktası kümeleri arasında kararlıydı; mekanizma atfımız değildi. Mekanizma iddiaları daha çok eğitim koşusu taşımalıdır.
\item \textbf{Katman bazlı analizlerde agregasyonu belirtin.} Aynı blokların CLS, yama ortalaması ve tüm jeton özetleri farklı kayma büyüklükleri ve farklı minimumlar verir.
\item \textbf{Dikkat ve lokalite iddialarını gerçek bir örneklem boyutunda sınayın.} Raporladığımız iki sıfır sonuç da küçük bir figür partisinde pozitif görünürdü.
\end{enumerate}

Bu kıyaslamada eşleşmiş standart çekişmeli eğitim altında iki aile arasında seçim yapan uygulayıcılar için: CNN daha yüksek mutlak gürbüz doğruluk vermekte ve üstünlük ölçtüğümüz her görünümde korunmaktadır. \%62{,}76'daki WideResNet referansı, tarif, kapasite ve ek verinin birlikte~\cite{gowal2020uncovering} gürbüzlüğü mimari seçiminden daha ileri taşıdığını göstermektedir; gürbüzlük stratejisi olarak mimari seçimine yatırım yapmadan önce tartılmaya değer.

% ----------------------------------------------------------------------------
% 5.3 Sinirliliklar
% ----------------------------------------------------------------------------
\subsection{Sınırlılıklar}
\label{subsec:limitations}

Çalışmamız bilinçli olarak dardır ve bu darlık sonuçları belirli biçimlerde sınırlar.

Protokol duyarlılığı sonucu üç veri kümesi ve tek model çifti üzerinde ölçülmüştür. CIFAR-100'de yayılım CIFAR-10'dakinden büyüktür ($13{,}58\pm1{,}71$'e karşı $10{,}45\pm0{,}76$ puan) ve işaret on iki ölçümün tamamında korunmaktadır (Bölüm~\ref{subsec:cifar100}); bu, CIFAR'a özgü bir şeyden değil hedefler arasındaki temiz doğruluk farkından geçen bir mekanizmayla tutarlıdır. Ancak bu açıklama eksiktir ve sloganı değil sınırı yazıyoruz: dört protokolün üçü için geçerlidir; başarılı kaynak protokolü ise \emph{kaynağın} gürbüzlüğüne bağlı ikinci bir sürücü ekler ve bu sürücü genel ilişkiyi tersine çevirecek kadar güçlüdür (Bölüm~\ref{subsec:transfer}). Her iki bağımlılığın işlevsel biçimini de güvenle belirtemiyoruz: kestirimler altı tohum-ve-veri-kümesi kümesine dayanmaktadır ve havuzumuzdaki küçük farklı çiftlerin tamamı harici WideResNet'i içerdiğinden, küçük temiz doğruluk farkı ile farklı eğitim tarifi burada birbirinden ayrılamamaktadır. Değişen temiz doğruluk farklarına sahip birden çok mimari çiftini tarayan bir çalışma, bu ölçümleri bir kalibrasyon eğrisine dönüştürecektir.

İki kontrol noktası kümesi arasındaki farkların nedenini yalıtamıyoruz. Düzeltilmiş protokol, doğrulama düzenini değiştirmiş ve aynı anda yeni bir eğitim koşusu üretmiştir. Veri miktarı açıklaması dışlanmıştır (düzeltilmiş temiz modeller daha az veriyle daha iyidir); ancak seçim sızıntısını olağan koşudan koşuya değişimden ayırmak, özdeş veriyle sızıntılı ve sızıntısız seçim altında eğiten bir ablasyon gerektirir. Karşıtlığı sızıntının ölçümü olarak değil, tek koşuluk mekanizma iddialarının kırılganlığına dair kanıt olarak raporluyoruz.

Bu makaledeki her ölçüm tek bir tehdit modelinin içinde yapılmıştır: $\eps = 8/255$ ile $\Linf$. Protokol yayılımının norm değişince ayakta kalıp kalmadığı ayrı bir sorudur ve kaydedilmiş ama henüz yanıtlanmamıştır. Böyle bir sınamanın neyi gösterip neyi gösteremeyeceğini belirtmek gerekir: modellerimiz $\Linf$ ile eğitilmiştir, dolayısıyla onları bir $L_2$ bütçesi altında değerlendirmek $L_2$ gürbüzlüğünü değil ölçüm eksenini yoklar ve çıkan mutlak değerler $L_2$ ile eğitilmiş referanslarla karşılaştırılamaz.

Mimari başına üç koşu tohum varyansını sınırlar ama tarif uzayını taramaz; küçük veri kümelerinde ViT çekişmeli eğitiminin tarife duyarlı olduğu bilinmektedir~\cite{mo2022adversarial, debenedetti2023light} ve difüzyonla üretilmiş sentetik veriyle ViT'lerin CIFAR-10'da WideResNet'leri geçtiği bildirilmiştir~\cite{wang2023better, wu2025vision}. ViT sonuçlarımız ulaşılabilir en iyi ViT gürbüzlüğünü değil, temel çekişmeli eğitimi betimler. ViT-Tiny hattı $32\times32$ girdileri model içinde $224\times224$'e büyütür; CIFAR'a özgü kontrolümüz gradyan sıralamasının büyütme artefaktı olmadığını gösterir, ama bu kontrol kapasite ve yama boyutunda farklıdır; dolayısıyla karıştırıcıyı kalibre eder, ikinci bir eşleşmiş çift sağlamaz. İki model parametre sayısında da ayrışır (11,2M'ye karşı 5,7M); mutlak farkın mimariye atfı bu yüzden kısmi kalır. Son olarak tek bütçeli bir $L_\infty$ tehdit modeli değerlendiriyoruz ve TGR sonucumuz, yöntemin yayımlanmış kurulumunun yeniden üretimi değil bir uyarlamasıdır.

% ----------------------------------------------------------------------------
% 5.4 Gelecek Calismalar
% ----------------------------------------------------------------------------
\subsection{Gelecek Yönelimler}
\label{subsec:future}

En doğrudan uzantı, protokol duyarlılığı ölçümünü genel bir ölçüme dönüştürmektir: kontrollü temiz doğruluk farklarına sahip mimari çiftlerini tarayıp protokoller arası yayılımın bu farkla nasıl ölçeklendiğini ölçmek. İlişki, üç hedefli korelasyonumuzun önerdiği kadar düzenliyse, yayımlanmış ham oran asimetrilerinin atılmak yerine yeniden yorumlanmasına imkân veren bir düzeltme sağlar.

Daha önce gelecek çalışma diye sıraladığımız uzantılardan ikisi bu arada koşulmuştur; hem neyi karara bağladıklarını hem neyi bağlamadıklarını kaydediyoruz. Sızıntı ablasyonu altı tam bütçeli yörünge üzerinde yürütülmüş ve sızıntının kontrol noktası seçimine etkisini \emph{saptamamıştır}; başlangıçta konan sorunun daha zayıf bir sürümünü yanıtlar, çünkü tedavi tek bir doğrulama bölmesine uygulanmıştır ve o tasarımda bölme kimliği ile sızıntı birbirinden ayrıştırılamaz. Bilgi verici yan ürünü, Bölüm~\ref{subsec:statistics_results}'te raporlanan seçim protokolü yayılımıdır. İkinci veri kümesi de koşulmuştur (Bölüm~\ref{subsec:cifar100}). Açık kalan, kapasitesi eşlenmiş çifttir (örneğin ViT-Small'a karşı ResNet-50); bu, ayakta kalan davranışsal farkların --- seyreklik sıralaması ve derinlik profilleri --- ailelerin mi yoksa bu küçük modellerin mi özelliği olduğunu test edecektir.

% ============================================================================
% BOLUM 6: SONUC
% ============================================================================
\section{Sonuç}
\label{sec:conclusion}

CNN--Görü Dönüştürücüsü gürbüzlük karşılaştırmalarındaki uyuşmazlığın ne kadarının modellerden değil ölçümden geldiğini sorduk ve bunu; modelleri, veriyi, tehdit modelini ve saldırı bütçesini sabit tutup yalnızca karşılaştırmanın nasıl puanlandığını değiştirerek yanıtladık.

Ölçüm hükmediyor. Özdeş çekişmeli eğitilmiş ResNet-18 ve ViT-Tiny çiftlerinde, mimariler arası transfer asimetrisi dört yerleşik koşullama protokolü boyunca $+4{,}4$ ile $+14{,}6$ puan arasında değişmektedir --- 3,3 katlık bir yayılım; aynı nicelik üzerinde ölçüldüğünde bu protokol etkisi, bağımsız eğitim koşuları arasındaki standart sapmanın $3{,}3$ ile $22{,}7$ katı arasındadır --- ve protokol, bir eşdeğerlik testinin kabul mü ret mi edileceğine karar vermektedir. Koşulsuz oranlar böyle bir karşılaştırma için hiçbir biçimde geçerli temel değildir: WideResNet-28-10 referansının eklendiği $3\times3$ matriste, koşullu oranlardan sapmaları hedefin temiz hatasıyla $r = 0{,}997$ ile açıklanmaktadır. Protokolün değiştirmediği şey yöndür; yön protokoller, tohumlar ve saldırı aileleri boyunca kararlıdır ve üç hedefli analizimiz onu CNN'de üretilen pertürbasyonların özel bir gücünden çok hedefin kendi kırılganlığıyla ($r = 0{,}986$) ilişkilendirmektedir.

Koşullu ayrıştırma, mutlak gürbüzlüğü temiz doğruluk çarpanı ile koşullu duyarlılık çarpanına ayırır. Düzeltilmiş koşularımızda ikisi de CNN lehinedir; ancak mutlak sıralaması özdeş olan önceki bir kontrol noktası kümesinde iki çarpan arasındaki pay farklıydı. Bu, mekanizma düzeyindeki iddiaların sıralama düzeyindeki iddialardan daha kırılgan olduğunu ve desteklenmek için daha çok eğitim koşusu gerektirdiğini gösterir.

Bu literatürde yaygın üç iddia doğrudan ölçümden sağ çıkamamıştır: çekişmeli eğitilmiş CNN gradyanları daha seyrektir ama mekânsal olarak daha lokalize değildir; CLS dikkat entropisi $n=1000$'de $L_\infty$ saldırısı altında değişmemektedir; dönüştürücüye özgü bir transfer saldırısı bu rejimde transfer kazancı sağlamamaktadır. İki davranışsal fark ise bağımsız kontrol noktası kümelerinde tekrarlanarak ayakta kalmıştır: mimarinin değil çekişmeli eğitimin yarattığı gradyan seyrekliği sıralaması ile çekişmeli hasarın derinlik profili --- CNN'de geç, yerel ve büyüklük değiştiren bir olay; ViT'te dağıtık, yalnızca yön değiştiren bir dönme.

Bu ölçümlerden, ilk ikisi en önemli olan altı raporlama gerekliliği damıttık: koşullama protokolünü belirtin ve eşleştirilmiş olanları yeğleyin; koşulsuz transfer oranlarını mimari karşılaştırması olarak sunmayın. Çalışmamız tek veri kümesini ve tek model çiftini kapsamaktadır; dolayısıyla protokol etkisinin diğer kurulumlardaki boyutu kalibre edilmeyi beklemektedir. Etkinin arkasındaki mekanizma --- hedefler arasındaki eşitsiz temiz doğruluk farkı --- ise kurulumumuza özgü değildir. Protokol karşılaştırmasının yeniden üretilip genişletilebilmesi için tam boru hattını, örnek bazında kayıtları ve analiz betiklerini yayımlıyoruz.


% ----------------------------------------------------------------------------
% VERI VE KOD ERISILEBILIRLIGI
% ----------------------------------------------------------------------------
\section*{Veri ve Kod Erişilebilirliği}
Bu çalışmadaki deneyler, \url{https://www.cs.toronto.edu/~kriz/cifar.html} adresinden herkese açık olan CIFAR-10 veri kümesini~\cite{krizhevsky2009learning} kullanmaktadır. Önceden eğitilmiş WideResNet-28-10 ağırlıkları RobustBench kıyaslamasından~\cite{croce2021robustbench, gowal2020uncovering} alınmıştır. Analiz hattının kaynak kodu, eğitilmiş model kontrol noktaları, örnek bazında değerlendirme kayıtları ve tüm analiz betikleri kabul sonrasında herkese açık hâle getirilecektir.

% ----------------------------------------------------------------------------
% TESEKKUR
% ----------------------------------------------------------------------------
\section*{Teşekkür}
Yazarlar, hesaplama kaynakları için Fırat Üniversitesi Yazılım Mühendisliği Bölümü'ne teşekkür eder.

% ----------------------------------------------------------------------------
% KAYNAKLAR
% ----------------------------------------------------------------------------
\balance
\bibliographystyle{IEEEtran}
\bibliography{../manuscript/references}

\end{document}
